\documentclass[11pt,a4paper]{article}
\usepackage[utf8]{inputenc}
\usepackage[T1]{fontenc}
\usepackage[spanish,es-minimal,es-noshorthands]{babel}
\usepackage{lmodern}
\linespread{1.22}
\usepackage[a4paper,top=2.4cm,bottom=2.2cm,left=1.5cm,right=1.5cm,headheight=1.1cm,headsep=0.5cm,footskip=1cm]{geometry}
\usepackage{amsmath,amssymb,cancel}
\usepackage{icomma}
\usepackage{graphicx}
\usepackage[table]{xcolor}
\usepackage{tikz}\usetikzlibrary{calc}
\usepackage{pgfplots}\pgfplotsset{compat=1.17}
\usepackage{fontawesome5}
\usepackage{ccicons}
\usepackage{fancyhdr}
\usepackage{enumitem}
\usepackage{booktabs}
\usepackage{tabularx}
\usepackage{needspace}
\usepackage[most]{tcolorbox}
\usepackage{titlesec}
\usepackage[hidelinks]{hyperref}
\hypersetup{pdftitle={Apuntes de teoría. Unidad 9: Combinatoria y probabilidad. 1.º de Bachillerato CCSS},pdfsubject={Matemáticas Aplicadas a las Ciencias Sociales I},pdfkeywords={principio de multiplicación, variaciones, permutaciones, combinaciones, números combinatorios, experimentos aleatorios, sucesos, regla de Laplace, probabilidad condicionada, probabilidad total, teorema de Bayes}}

% colors de la web (FentMates, bloque de Probabilidad y Estadística: morado; es manté el nom fmgreen dels apunts de les altres unitats)
\definecolor{fmnavy}{HTML}{1B2748}
\definecolor{fmgreen}{HTML}{684E88}
\definecolor{fmgreenl}{HTML}{8064A1}
\definecolor{fmyellow}{HTML}{FFC94D}
\definecolor{fmgold}{HTML}{E9A919}
\definecolor{fmaccent}{HTML}{FF7A50}
\definecolor{fmteal}{HTML}{2BB3A3}
\definecolor{fmdef}{HTML}{1D4F73}\definecolor{fmdefbg}{HTML}{EEF4F9}
\definecolor{fmprop}{HTML}{3B4252}\definecolor{fmpropbg}{HTML}{F3F4F7}
\definecolor{fmex}{HTML}{A3203A}\definecolor{fmexbg}{HTML}{FBF0F2}
\definecolor{fmmuted}{HTML}{5D6685}\definecolor{fmball0}{HTML}{E2483D}\definecolor{fmball1}{HTML}{2F6FD6}
\definecolor{fmsurface}{HTML}{EEF1F8}
\definecolor{fmA}{HTML}{3B82F6}\definecolor{fmB}{HTML}{F97316}\definecolor{fmR}{HTML}{E11D48}\definecolor{fmK}{HTML}{16A34A}\definecolor{fmV}{HTML}{7C3AED}

% blocs que no es poden partir: es mesuren i es reserva exactament el seu espai (amb el títol o la capçalera del requadre)
\newsavebox{\fmkb}
\newcommand{\fmmeasure}[2]{\setbox\fmkb=\vbox{\hsize=\dimexpr\linewidth-#1\relax\linewidth=\hsize #2}}
\newcommand{\fmblock}[1]{\fmmeasure{0pt}{#1}\Needspace{\dimexpr\ht\fmkb+\dp\fmkb\relax}#1}
\newcommand{\fmheadblock}[2]{\fmmeasure{0pt}{#2}\Needspace{\dimexpr\ht\fmkb+\dp\fmkb+2.2cm\relax}#1#2}
\newcommand{\fmboxneed}[1]{\fmmeasure{9mm}{#1}\Needspace{\dimexpr\ht\fmkb+\dp\fmkb+1.6cm\relax}}
\newif\iffmstack\newif\iffmfit
\newcommand{\fmfits}[1]{\par\ifdim\dimexpr\pagegoal-\pagetotal\relax<\dimexpr\ht\fmkb+\dp\fmkb+#1\relax\global\fmfitfalse\else\global\fmfittrue\fi}
\newcommand{\fmtwo}[2]{\fmmeasure{0pt}{#1}\fmfits{0pt}\iffmfit #1\else #2\fi}
\newcommand{\fmheadtwo}[3]{\fmmeasure{0pt}{#2}\fmfits{2.2cm}\iffmfit #1#2\else\Needspace{\dimexpr1.6cm+7\baselineskip\relax}#1#3\fi}
\newcommand{\fmtwoin}[2]{\iffmstack #2\else #1\fi}
\newcommand{\fmboxneedalt}[1]{\fmmeasure{9mm}{#1}\fmfits{1.6cm}\iffmfit\global\fmstackfalse\else\global\fmstacktrue\Needspace{8\baselineskip}\fi}
\newcommand{\fmheadboxalt}[2]{\fmmeasure{9mm}{#2}\fmfits{3.4cm}\iffmfit\global\fmstackfalse\else\global\fmstacktrue\Needspace{\dimexpr2.4cm+5\baselineskip\relax}\fi#1}
\newcommand{\fmheadbox}[2]{\fmmeasure{9mm}{#2}\Needspace{\dimexpr\ht\fmkb+\dp\fmkb+3.4cm\relax}#1}
\newcommand{\lt}{<}\newcommand{\gt}{>}
% fórmula en línia que no es parteix mai; si no cap en una línia, es posa centrada
% fórmula en línia; si ocupa més del 80 % de la línia, va centrada (amb la puntuació que la segueix)
\makeatletter
\newcommand{\fmm}[1]{\setbox0=\hbox{$#1$}\ifdim\wd0>0.8\linewidth\def\fmmnext{\fmmdisp{#1}}\else\def\fmmnext{\mbox{\box0}}\fi\fmmnext}
\def\fmmdisp#1{\@ifnextchar.{\fmmdispp{#1}}{\@ifnextchar,{\fmmdispc{#1}}{\@ifnextchar;{\fmmdisps{#1}}{\fmdisp{#1}}}}}
\def\fmmdispp#1.{\fmdisp{#1\,.}}\def\fmmdispc#1,{\fmdisp{#1\,,}}\def\fmmdisps#1;{\fmdisp{#1\,;}}
\makeatother
% fórmula centrada; si és més ampla que la línia, s'encongeix fins que hi cap
\newcommand{\fmdisp}[1]{\[\setbox2=\hbox{$\displaystyle #1$}\ifdim\wd2>\linewidth\resizebox{\linewidth}{!}{\box2}\else\box2\fi\]}
\DeclareUnicodeCharacter{2212}{-}
\DeclareUnicodeCharacter{2225}{\ensuremath{\parallel}}\DeclareUnicodeCharacter{22A5}{\ensuremath{\perp}}\DeclareUnicodeCharacter{2226}{\ensuremath{\nparallel}}
\DeclareUnicodeCharacter{207A}{\ensuremath{^{+}}}\DeclareUnicodeCharacter{207B}{\ensuremath{^{-}}}
\setlength{\parindent}{0pt}
\setlength{\parskip}{12pt}
\setlength{\lineskip}{5pt}\setlength{\lineskiplimit}{5pt}  % més aire al voltant de les fraccions en línia
\setlength{\emergencystretch}{3em}
\setlist{itemsep=4pt,topsep=4pt}
\renewcommand{\thesection}{9.\arabic{section}}
\titleformat{\section}{\color{fmnavy}\Large\bfseries}{\colorbox{fmdef}{\color{white}\strut\ \thesection\ }}{0.6em}{}[{\color{fmdef}\titlerule[1.2pt]}]
\titlespacing*{\section}{0pt}{10pt}{4pt}
\titleformat*{\subsection}{\color{fmgreen}\large\bfseries}
\titlespacing*{\subsection}{0pt}{6pt plus 2pt}{0pt}
\allowdisplaybreaks

% dibuix de la targeta de la unitat a la web (illustrations.js), en blanc
\newcommand{\fmcard}{\begin{tikzpicture}[x=0.062cm,y=-0.062cm,white,line width=3pt,line cap=round,line join=round]
  \begin{scope}[transparency group,opacity=0.15]
  \fill (20,60) circle[radius=5];
  \draw (20,60)--(80,26) (20,60)--(80,94) (80,26)--(130,12) (80,26)--(130,40) (80,94)--(130,80) (80,94)--(130,108);
  \fill (80,26) circle[radius=4] (80,94) circle[radius=4] (130,12) circle[radius=3.5] (130,40) circle[radius=3.5] (130,80) circle[radius=3.5] (130,108) circle[radius=3.5];
  \node[anchor=base west,inner sep=0pt,font=\fontsize{18}{18}\selectfont\itshape] at (36,36) {0,6};
  \node[anchor=base west,inner sep=0pt,font=\fontsize{18}{18}\selectfont\itshape] at (36,92) {0,4};
  \begin{scope}[rotate around={-12:(190,58)}]
    \draw[rounded corners=0.74cm] (162,30) rectangle (218,86);
    \fill (176,44) circle[radius=4.5] (190,58) circle[radius=4.5] (204,72) circle[radius=4.5];
  \end{scope}
  \end{scope}
\end{tikzpicture}}
\newcommand{\fmbrand}{\raisebox{-0.25\height}{\includegraphics[height=0.5cm]{logo-fentmates.pdf}}}

% capçalera i peu
\pagestyle{fancy}\fancyhf{}
\fancyhead[L]{\small\fmbrand}
\fancyhead[R]{\small\color{fmmuted}1.º Bachillerato CCSS · Unidad 9: Combinatoria y probabilidad}
\fancyfoot[L]{\scriptsize\color{fmmuted}\href{https://creativecommons.org/licenses/by-nc-sa/4.0/deed.ca}{\ccbyncsa}\ \ FentMates · fentmates.org · CC BY-NC-SA 4.0}
\fancyfoot[R]{\small\color{fmmuted}\thepage}
\renewcommand{\headrulewidth}{0.4pt}\renewcommand{\headrule}{\color{fmmuted!40}\hrule height \headrulewidth}
\fancypagestyle{primera}{\fancyhf{}\renewcommand{\headrulewidth}{0pt}
  \fancyfoot[L]{\scriptsize\color{fmmuted}\href{https://creativecommons.org/licenses/by-nc-sa/4.0/deed.ca}{\ccbyncsa}\ \ FentMates · fentmates.org · CC BY-NC-SA 4.0}
  \fancyfoot[R]{\small\color{fmmuted}\thepage}}

% caixes (mateixos colors que la web)
% quan una caixa es parteix entre dues pàgines, la vora del tall és discontínua
\tcbset{fmbase/.style={enhanced,breakable,arc=3mm,boxrule=1pt,left=3.5mm,right=3.5mm,top=2mm,bottom=2.5mm,before skip=10pt,after skip=10pt,fonttitle=\bfseries,
  frame style={fill=tcbcolframe},
  overlay first={\fill[tcbcolback] ([xshift=1pt]frame.south west) rectangle ([xshift=-1pt,yshift=1.2pt]frame.south east);
    \draw[tcbcolframe,line width=1pt,dash pattern=on 5pt off 3pt] (frame.south west)--(frame.south east);},
  overlay middle={\fill[tcbcolback] ([xshift=1pt]frame.south west) rectangle ([xshift=-1pt,yshift=1.2pt]frame.south east);
    \fill[tcbcolback] ([xshift=1pt]frame.north west) rectangle ([xshift=-1pt,yshift=-1.2pt]frame.north east);
    \draw[tcbcolframe,line width=1pt,dash pattern=on 5pt off 3pt] (frame.south west)--(frame.south east) (frame.north west)--(frame.north east);},
  overlay last={\fill[tcbcolback] ([xshift=1pt]frame.north west) rectangle ([xshift=-1pt,yshift=-1.2pt]frame.north east);
    \draw[tcbcolframe,line width=1pt,dash pattern=on 5pt off 3pt] (frame.north west)--(frame.north east);}}}
\newtcolorbox{fmdef}[1]{fmbase,colback=fmdefbg,colframe=fmdef,coltitle=white,title={#1}}
\newtcolorbox{fmprop}[1]{fmbase,colback=fmpropbg,colframe=fmprop,coltitle=white,title={#1},halign title=center}
\newtcolorbox{fmex}[1]{fmbase,before upper=\raggedright,colback=fmexbg,colframe=fmex,colbacktitle=fmexbg,coltitle=fmex,title={\faPencil*\ #1},
  titlerule=0.4pt,titlerule style={fmex!40,dashed}}
\newtcolorbox{fmwarn}[1]{fmbase,colback=fmsf!6,colframe=fmsf,coltitle=white,title={#1}}
\newtcolorbox{fmnote}[1][]{enhanced,breakable,colback=white,colframe=white,borderline west={3pt}{0pt}{fmaccent},arc=0pt,left=3mm,right=2mm,top=1mm,bottom=1mm,before skip=6pt,after skip=6pt,before upper=\raggedright,#1}
\newcommand{\fmkey}[1]{\colorbox{fmyellow!45}{\strut #1}}
\newcommand{\fmuline}[1]{\textbf{\color{fmaccent}#1}}
\definecolor{fmok}{HTML}{15803D}\definecolor{fmev}{HTML}{2563EB}\definecolor{fmsf}{HTML}{C2410C}\definecolor{fmsi}{HTML}{7C3AED}\definecolor{fmas}{HTML}{A3203A}
\newcommand{\fmbadge}[2]{\tcbox[on line,colback=#1,colframe=#1,arc=2mm,boxsep=0pt,left=4pt,right=4pt,top=1.5pt,bottom=1.5pt]{\color{white}\footnotesize\bfseries #2}}
\newcommand{\fmchip}[1]{\tcbox[on line,colback=fmsf!10,colframe=fmsf!50,arc=2mm,boxsep=1pt,left=4pt,right=4pt]{#1}}
\newcommand{\fmlevel}[2]{\tcbox[on line,colback=white,colframe=fmmuted!40,arc=2mm,boxsep=2pt,halign=center]{\begin{tabular}{c}{\scriptsize\color{fmmuted}\MakeUppercase{#1}}\\#2\end{tabular}}}
% gràfiques
\pgfplotsset{fmaxis/.style={axis lines=middle,grid=major,grid style={fmmuted!15},axis line style={fmmuted,-latex},
  tick label style={font=\scriptsize,color=fmmuted},clip=true,unbounded coords=jump,restrict y to domain=-200:200,
  every axis plot/.append style={line join=round}},
  fmcurve/.style={fmaccent,line width=1.6pt,smooth}}
\tikzset{fmasy/.style={fmmuted,dashed,line width=0.9pt},fmguide/.style={fmaccent,dotted,line width=0.8pt},
  fmdot/.style={circle,fill=fmaccent,inner sep=0pt,minimum size=5pt},fmopen/.style={circle,draw=fmaccent,fill=white,line width=1pt,inner sep=0pt,minimum size=5pt}}
\newcommand{\fmpau}[1]{\tcbox[on line,colback=fmnavy,colframe=fmnavy,arc=2mm,boxsep=0pt,left=4pt,right=4pt,top=1.5pt,bottom=1.5pt]{\color{white}\scriptsize\bfseries\MakeUppercase{#1}}\ }

\begin{document}
\thispagestyle{primera}
\begin{tikzpicture}[remember picture,overlay]
  \fill[fmgreen] (current page.north west) rectangle ([yshift=-10.1cm]current page.north east);
  \begin{scope}
    \clip (current page.north west) rectangle ([yshift=-10.1cm]current page.north east);
    \node[anchor=north west,inner sep=0pt] at ([xshift=6.6cm,yshift=-1.3cm]current page.north west) {\fmcard};
  \end{scope}
  \node[anchor=north west,inner sep=0pt] at ([xshift=1.5cm,yshift=-1.1cm]current page.north west) {\tcbox[colback=white,colframe=white,arc=2.5mm,boxsep=0pt,left=7pt,right=7pt,top=5pt,bottom=5pt]{\includegraphics[height=0.75cm]{logo-fentmates.pdf}}};
\end{tikzpicture}
\vspace*{0.55cm}
\tcbox[colback=white!15!fmgreen,colframe=white!15!fmgreen,arc=2.5mm,boxsep=0pt,left=6pt,right=6pt,top=3pt,bottom=3pt]{\color{white}\small\bfseries Apuntes de teoría · 1.º Bachillerato · Matemáticas Aplicadas a las Ciencias Sociales I}

\vspace{4pt}
{\raggedright\color{white}\fontsize{25}{30}\selectfont\bfseries{\color{fmyellow}Unidad 9} · Combinatoria y probabilidad\par}
\vspace{6pt}
\begin{minipage}{\linewidth}\color{white!88!fmgreen}\small Teoría y ejemplos resueltos paso a paso: técnicas de recuento (diagramas de árbol, principio de multiplicación, factorial), combinatoria (variaciones, permutaciones y combinaciones, con repetición y sin ella, y números combinatorios), experimentos aleatorios y sucesos, la probabilidad y la regla de Laplace, el cálculo de probabilidades con sus fórmulas y con tablas de contingencia, la probabilidad condicionada, los sucesos dependientes e independientes, las experiencias compuestas (árboles, extracciones con reemplazamiento y sin él) y los teoremas de la probabilidad total y de Bayes.\end{minipage}\par
% l'índex comença sempre a la mateixa altura, per davall de la franja (com en les unitats 5-10)
\vspace*{\dimexpr 8.6cm-\pagetotal\relax}

\begin{tcolorbox}[colback=fmsurface,colframe=fmsurface,arc=3mm,left=4mm,right=4mm,top=2mm,bottom=2mm]
\small\textbf{Índice}\par\smallskip
\begin{minipage}[t]{0.49\linewidth}\small \mbox{\hyperlink{sec1}{\textbf{\color{fmgreen}9.1}\ \ Técnicas de recuento}}\\[3pt]
\mbox{\hyperlink{sec2}{\textbf{\color{fmgreen}9.2}\ \ Variaciones}}\\[3pt]
\mbox{\hyperlink{sec3}{\textbf{\color{fmgreen}9.3}\ \ Permutaciones}}\\[3pt]
\mbox{\hyperlink{sec4}{\textbf{\color{fmgreen}9.4}\ \ Combinaciones. Números combinatorios}}\\[3pt]
\mbox{\hyperlink{sec5}{\textbf{\color{fmgreen}9.5}\ \ Espacio muestral y sucesos}}\\[3pt]
\mbox{\hyperlink{sec6}{\textbf{\color{fmgreen}9.6}\ \ Operaciones con sucesos}}\\[3pt]
\mbox{\hyperlink{sec7}{\textbf{\color{fmgreen}9.7}\ \ La probabilidad}}\\[3pt]\end{minipage}\hfill\begin{minipage}[t]{0.49\linewidth}\small \mbox{\hyperlink{sec8}{\textbf{\color{fmgreen}9.8}\ \ Regla de Laplace}}\\[3pt]
\mbox{\hyperlink{sec9}{\textbf{\color{fmgreen}9.9}\ \ Fórmulas de la probabilidad}}\\[3pt]
\hyperlink{sec10}{\textbf{\color{fmgreen}9.10}\ \ Cálculo de probabilidades: con el espacio muestral y sin él}\\[3pt]
\mbox{\hyperlink{sec11}{\textbf{\color{fmgreen}9.11}\ \ Probabilidad condicionada}}\\[3pt]
\mbox{\hyperlink{sec12}{\textbf{\color{fmgreen}9.12}\ \ Sucesos dependientes e independientes}}\\[3pt]
\mbox{\hyperlink{sec13}{\textbf{\color{fmgreen}9.13}\ \ Experiencias compuestas}}\\[3pt]
\mbox{\hyperlink{sec14}{\textbf{\color{fmgreen}9.14}\ \ Probabilidad total y teorema de Bayes}}\\[3pt]\end{minipage}
\end{tcolorbox}

\Needspace{6.2cm}
\section{Técnicas de recuento. Diagramas de árbol, principio de multiplicación y factorial}\hypertarget{sec1}{}

Para calcular probabilidades a menudo tenemos que \textbf{contar} cuántos casos posibles y cuántos casos favorables hay. Cuando son pocos, los podemos escribir todos; cuando son muchos (¿cuántos números de teléfono hay? ¿cuántas apuestas de la lotería?), necesitamos técnicas para contarlos sin tener que escribirlos. Estas técnicas se llaman \textbf{técnicas de recuento} o \textbf{combinatoria}.

\Needspace{5\baselineskip}
\begin{fmdef}{Diagrama de árbol y principio de multiplicación}
Cuando una elección se hace en varias \textbf{etapas} sucesivas, podemos representar todas las posibilidades con un \textbf{diagrama de árbol}: de cada punto sale una rama para cada opción de la etapa siguiente. Cada camino desde la raíz hasta el final es un resultado.

\textbf{Principio de multiplicación}: si la primera etapa se puede hacer de \fmm{n_1} maneras, la segunda de \fmm{n_2} maneras, \dots{}, y la última de \fmm{n_k} maneras, el número total de posibilidades es

\fmdisp{n_1\cdot n_2\cdots n_k}

\end{fmdef}

\Needspace{5\baselineskip}
\begin{fmex}{Ejemplos resueltos}
\begin{enumerate}[label=\textbf{\arabic*.},leftmargin=*]
\item Un restaurante ofrece un menú con 3 primeros platos, 4 segundos y 2 postres. ¿Cuántos menús distintos se pueden elegir?

Hay tres etapas: elegir el primero (3 opciones), el segundo (4) y el postre (2). En el árbol, de la raíz salen 3 ramas, de cada una 4 y de cada una de estas, 2.

\fmm{3\cdot4\cdot2=24} menús distintos.
\item Las matrículas de los coches tienen 4 cifras y 3 letras (se usan 20 letras). ¿Cuántas matrículas distintas puede haber?

Cada cifra se puede elegir de 10 maneras y cada letra de 20 (se pueden repetir):

\fmm{10\cdot10\cdot10\cdot10\cdot20\cdot20\cdot20=10^4\cdot20^3=80\,000\,000} matrículas.
\item Con las cifras 1, 2, 3, 4 y 5, ¿cuántos números de tres cifras se pueden escribir (a) si se pueden repetir las cifras, (b) si no se pueden repetir, (c) que sean pares y sin repetir cifras?

(a) Cada cifra se puede elegir de 5 maneras: \fmm{5\cdot5\cdot5=125}.

(b) La primera, de 5 maneras; la segunda, de 4 (no puede ser la primera); la tercera, de 3: \fmm{5\cdot4\cdot3=60}.

(c) Empezamos por la cifra con condición: las unidades tienen que ser 2 o 4 (2 maneras). Después, las centenas (4 maneras) y las decenas (3 maneras): \fmm{2\cdot4\cdot3=24}.
\end{enumerate}
\end{fmex}

\Needspace{5\baselineskip}
\begin{fmprop}{Principio de la suma. Contar por el contrario}
Si los casos se pueden separar en grupos \textbf{que no tienen ningún elemento en común}, el total es la \textbf{suma} de los casos de cada grupo. Y, a menudo, es más fácil contar los casos que \textbf{no} queremos y restarlos del total.

\textbf{Ejemplo}: ¿cuántas contraseñas de 4 cifras tienen al menos una cifra repetida? Todas: \fmm{10^4=10\,000}; sin ninguna repetida: \fmm{10\cdot9\cdot8\cdot7=5\,040}. Con alguna repetida: \fmm{10\,000-5\,040=4\,960}.

\textbf{Ejemplo}: ¿cuántas matrículas de motocicleta hay si pueden tener 2 o 3 letras (de 20) seguidas de 2 cifras? Con 2 letras: \fmm{20^2\cdot10^2=40\,000}; con 3: \fmm{20^3\cdot10^2=800\,000}. En total, \fmm{840\,000}.
\end{fmprop}

\Needspace{5\baselineskip}
\begin{fmdef}{El factorial de un número}
En muchos recuentos aparecen productos de números naturales consecutivos. Para escribirlos de forma corta usamos el \textbf{factorial}: si \fmm{n} es un número natural,

\fmdisp{n!=n\cdot(n-1)\cdot(n-2)\cdots2\cdot1\qquad\qquad 0!=1}

Por ejemplo, \fmm{4!=4\cdot3\cdot2\cdot1=24} y \fmm{5!=120}. Observa que \fmm{n!=n\cdot(n-1)!}, lo que permite simplificar cocientes: \fmm{\dfrac{7!}{5!}=\dfrac{7\cdot6\cdot5!}{5!}=7\cdot6=42}. En la calculadora es la función \textbf{x!}

\begin{fmnote}Los factoriales crecen muy deprisa: \fmm{10!=3\,628\,800} y \fmm{20!} ya tiene 19 cifras.\end{fmnote}
\end{fmdef}

\Needspace{8\baselineskip}
\section{Variaciones con repetición y sin repetición}\hypertarget{sec2}{}

Muchos problemas de recuento consisten en elegir \fmm{k} elementos de un conjunto de \fmm{n} elementos y \textbf{colocarlos en un orden} (primero, segundo, tercero\dots{}). Según si los elementos se pueden repetir o no, hablamos de variaciones con repetición o sin repetición.

\Needspace{5\baselineskip}
\begin{fmdef}{Variaciones}
\begin{itemize}
\item \textbf{Variaciones sin repetición} de \fmm{n} elementos tomados de \fmm{k} en \fmm{k}: grupos ordenados de \fmm{k} elementos \textbf{distintos}. Importa el orden (AB $\neq$ BA). El primer puesto tiene \fmm{n} opciones, el segundo \fmm{n-1}\dots{}

\fmdisp{V_{n,k}=\underbrace{n\cdot(n-1)\cdot(n-2)\cdots}_{k\ \text{factores}}=\frac{n!}{(n-k)!}}

\item \textbf{Variaciones con repetición}: grupos ordenados de \fmm{k} elementos que \textbf{se pueden repetir}. Cada puesto tiene \fmm{n} opciones:

\fmdisp{VR_{n,k}=n\cdot n\cdots n=n^k}

\end{itemize}
\end{fmdef}

\Needspace{5\baselineskip}
\begin{fmdef}{En la calculadora}
\begin{fmnote}\fmm{V_{n,k}} es la función \textbf{nPr}; en la fx-991SP CW, la encontrarás (con \textbf{x!} y \textbf{nCr}) en el menú CATÁLOGO $\to$ Probabilidad. Por ejemplo, \fmm{V_{8,3}}: 8 nPr 3 = 336. Para \fmm{VR_{n,k}=n^k} basta con la potencia.\end{fmnote}
\end{fmdef}

\Needspace{5\baselineskip}
\begin{fmex}{Ejemplos resueltos}
\begin{enumerate}[label=\textbf{\arabic*.},leftmargin=*]
\item En una carrera participan 8 corredores. ¿De cuántas maneras pueden quedar repartidas las medallas de oro, plata y bronce?

Elegimos 3 corredores de 8, \textbf{importa el orden} (no es lo mismo ganar el oro que la plata) y no se repiten (nadie gana dos medallas): variaciones sin repetición.

\fmm{V_{8,3}=8\cdot7\cdot6=336}
\item Una quiniela de fútbol tiene 14 partidos y en cada uno se marca 1, X o 2. ¿Cuántas quinielas distintas se pueden rellenar?

Para cada partido elegimos uno de los 3 signos, importa el orden (el partido en el que lo ponemos) y se pueden repetir: variaciones con repetición.

\fmm{VR_{3,14}=3^{14}=4\,782\,969}
\item Con las cifras 0, 1, 2, 3 y 4, ¿cuántos números de tres cifras distintas se pueden formar? ¿Cuántos son múltiplos de 5?

Todas las ordenaciones de 3 cifras distintas son \fmm{V_{5,3}=60}, pero las que empiezan por 0 no son de tres cifras: hay \fmm{V_{4,2}=12}. Quedan \fmm{60-12=48}. (O bien: la primera, de 4 maneras; después, \fmm{V_{4,2}=12}: \fmm{4\cdot12=48}.)

Múltiplos de 5: tienen que acabar en 0 (no hay ningún 5). Fijamos el 0 al final y elegimos las otras dos cifras ordenadas entre 1, 2, 3 y 4: \fmm{V_{4,2}=12}.
\end{enumerate}
\end{fmex}

\Needspace{7.3cm}
\section{Permutaciones ordinarias, con repetición y circulares}\hypertarget{sec3}{}
\begin{fmdef}{Permutaciones}
Las \textbf{permutaciones} de \fmm{n} elementos son las maneras de \textbf{ordenar todos} los \fmm{n} elementos. Son variaciones de \fmm{n} elementos tomados de \fmm{n} en \fmm{n}:

\fmdisp{P_n=V_{n,n}=n\cdot(n-1)\cdots2\cdot1=n!}

\end{fmdef}

\Needspace{5\baselineskip}
\begin{fmdef}{Permutaciones con repetición}
Si queremos ordenar \fmm{n} elementos entre los que hay \fmm{a} iguales de un tipo, \fmm{b} iguales de otro\dots{}, (\fmm{a+b+\dots=n}), intercambiar entre sí dos elementos iguales no da una ordenación nueva. Cada ordenación aparece \fmm{a!\,b!\cdots} veces entre las \fmm{n!}, así que:

\fmdisp{PR_n^{\,a,b,\dots}=\frac{n!}{a!\;b!\cdots}}

\textbf{Ejemplo}: con las letras de CASA (dos A) se pueden hacer \fmm{PR_4^{\,2}=\frac{4!}{2!}=12} ordenaciones, y no 24: si en CASA intercambiamos las dos A, obtenemos la misma palabra.
\end{fmdef}

\Needspace{5\baselineskip}
\begin{fmdef}{Permutaciones circulares}
Si \fmm{n} personas se sientan alrededor de una \textbf{mesa redonda} y solo importa a quién tiene cada uno a cada lado (no la silla concreta), girar a todos una silla no cambia nada. Fijamos una persona y ordenamos a las demás:

\fmdisp{PC_n=(n-1)!}

\end{fmdef}

\Needspace{5\baselineskip}
\begin{fmex}{Ejemplos resueltos}
\begin{enumerate}[label=\textbf{\arabic*.},leftmargin=*]
\item Cinco amigos, entre los cuales están Ana y Biel, se sientan en un banco de 5 plazas. ¿De cuántas maneras se pueden sentar? ¿Y si Ana y Biel quieren estar juntos?

Ordenamos a todos los amigos: \fmm{P_5=5!=120} maneras.

Si van juntos, los «pegamos» como si fueran una sola persona: ordenamos 4 «personas», \fmm{4!=24}, y, además, Ana y Biel se pueden intercambiar entre ellos (2 maneras): \fmm{2\cdot24=48}.
\item ¿Cuántas ordenaciones distintas se pueden hacer con las letras A, B, C, D y E? ¿Cuántas empiezan por vocal?

\fmm{P_5=5!=120}

La primera letra tiene que ser A o E (2 opciones) y las otras 4 se ordenan de \fmm{4!} maneras: \fmm{2\cdot4!=48}.
\item ¿Cuántas palabras distintas (con sentido o sin él) se pueden formar con las letras de BANANA?

Hay 6 letras: 3 A, 2 N y 1 B. Son permutaciones con repetición: \fmm{PR_6^{\,3,2,1}=\dfrac{6!}{3!\,2!\,1!}=\dfrac{720}{12}=60}.
\item Una moneda se lanza 8 veces. ¿En cuántas de las secuencias posibles hay exactamente 3 caras?

Cada secuencia es una ordenación de 3 C y 5 X: \fmm{PR_8^{\,3,5}=\dfrac{8!}{3!\,5!}=56}. Es lo mismo que elegir qué 3 lanzamientos (de 8) son cara: \fmm{\binom83=56} (9.4).
\item ¿De cuántas maneras se pueden sentar 6 amigos alrededor de una mesa redonda? ¿Y si dos de ellos tienen que estar juntos?

\fmm{PC_6=5!=120}

Los «pegamos»: 5 «personas» en la mesa, \fmm{PC_5=4!=24}, por las 2 maneras de colocarse ellos dos: \fmm{48}.
\end{enumerate}
\end{fmex}

\Needspace{7.3cm}
\section{Combinaciones. Números combinatorios. Cómo reconocer cada tipo de recuento}\hypertarget{sec4}{}
\begin{fmdef}{Combinaciones y números combinatorios}
Las \textbf{combinaciones} de \fmm{n} elementos tomados de \fmm{k} en \fmm{k} son los grupos de \fmm{k} elementos distintos en los que \textbf{no importa el orden}: \{A, B\} y \{B, A\} son el mismo grupo.

Cada grupo de \fmm{k} elementos se puede ordenar de \fmm{k!} maneras, así que entre las variaciones cada grupo aparece \fmm{k!} veces. Por eso:

\fmdisp{C_{n,k}=\binom nk=\frac{V_{n,k}}{k!}=\frac{n!}{k!\,(n-k)!}}

El número \fmm{\binom nk} se lee «\fmm{n} sobre \fmm{k}» y se llama \textbf{número combinatorio}. En la calculadora es la función \textbf{nCr}.

\textbf{Ejemplo}: con las letras A, B, C y D, las variaciones de 2 son 12 (AB, BA, AC, CA\dots{}), pero los grupos (combinaciones) solo son \fmm{\frac{12}{2}=6}: \{A,B\}, \{A,C\}, \{A,D\}, \{B,C\}, \{B,D\}, \{C,D\}.
\end{fmdef}

\fmboxneed{\par\noindent \begin{minipage}[t]{0.48\linewidth}\vspace{0pt}\raggedright \begin{itemize}
\item \fmm{\binom n0=\binom nn=1}, \fmm{\ \binom n1=\binom n{n-1}=n}
\item \textbf{Complementarios}: \fmm{\binom nk=\binom n{n-k}} (elegir los \fmm{k} que entran es lo mismo que elegir los \fmm{n-k} que no entran).
\item \textbf{Fórmula de Stifel}: \fmm{\binom{n-1}{k-1}+\binom{n-1}{k}=\binom nk}. Cada número del triángulo es la suma de los dos que tiene encima.
\item La suma de la fila \fmm{n} es \fmm{\binom n0+\binom n1+\dots+\binom nn=2^n}: es el número total de subconjuntos de un conjunto de \fmm{n} elementos.
\end{itemize}

\textbf{Binomio de Newton}: los números combinatorios son los coeficientes de las potencias de un binomio,

\fmdisp{(a+b)^n=\binom n0a^n+\binom n1a^{n-1}b+\dots+\binom nn b^n}

Por ejemplo, \fmm{(a+b)^4=a^4+4a^3b+6a^2b^2+4ab^3+b^4} (fila 4 del triángulo).\end{minipage}\hfill\begin{minipage}[t]{0.48\linewidth}\vspace{0pt}\raggedright \begin{center}\renewcommand{\arraystretch}{1.5}\begin{tabular}{cc}\toprule
$n=0$ & 1\\
$n=1$ & 1   1\\
$n=2$ & 1   2   1\\
$n=3$ & 1   3   3   1\\
$n=4$ & 1   4   6   4   1\\
$n=5$ & 1   5   10   10   5   1\\
$n=6$ & 1   6   15   20   15   6   1\\
\bottomrule\end{tabular}\end{center}\end{minipage}\par\medskip}
\begin{fmprop}{Propiedades de los números combinatorios. Triángulo de Tartaglia}
\par\noindent \begin{minipage}[t]{0.48\linewidth}\vspace{0pt}\raggedright \begin{itemize}
\item \fmm{\binom n0=\binom nn=1}, \fmm{\ \binom n1=\binom n{n-1}=n}
\item \textbf{Complementarios}: \fmm{\binom nk=\binom n{n-k}} (elegir los \fmm{k} que entran es lo mismo que elegir los \fmm{n-k} que no entran).
\item \textbf{Fórmula de Stifel}: \fmm{\binom{n-1}{k-1}+\binom{n-1}{k}=\binom nk}. Cada número del triángulo es la suma de los dos que tiene encima.
\item La suma de la fila \fmm{n} es \fmm{\binom n0+\binom n1+\dots+\binom nn=2^n}: es el número total de subconjuntos de un conjunto de \fmm{n} elementos.
\end{itemize}

\textbf{Binomio de Newton}: los números combinatorios son los coeficientes de las potencias de un binomio,

\fmdisp{(a+b)^n=\binom n0a^n+\binom n1a^{n-1}b+\dots+\binom nn b^n}

Por ejemplo, \fmm{(a+b)^4=a^4+4a^3b+6a^2b^2+4ab^3+b^4} (fila 4 del triángulo).\end{minipage}\hfill\begin{minipage}[t]{0.48\linewidth}\vspace{0pt}\raggedright \begin{center}\renewcommand{\arraystretch}{1.5}\begin{tabular}{cc}\toprule
$n=0$ & 1\\
$n=1$ & 1   1\\
$n=2$ & 1   2   1\\
$n=3$ & 1   3   3   1\\
$n=4$ & 1   4   6   4   1\\
$n=5$ & 1   5   10   10   5   1\\
$n=6$ & 1   6   15   20   15   6   1\\
\bottomrule\end{tabular}\end{center}\end{minipage}\par\medskip
\end{fmprop}

\Needspace{5\baselineskip}
\begin{fmprop}{¿Cómo sabemos qué tipo de recuento es?}
Cuando elegimos \fmm{k} elementos de un conjunto de \fmm{n}, hagámonos estas preguntas:

        \textbf{Importa el orden y se pueden repetir}Variaciones con repetición: \fmm{VR_{n,k}=n^k}. Códigos, matrículas, quinielas, lanzamientos repetidos\dots{}
        \textbf{Importa el orden, no se repiten y no entran todos}Variaciones sin repetición: \fmm{V_{n,k}=n(n-1)\cdots} (\fmm{k} factores). Podios, cargos (presidente, secretario\dots{}), banderas de colores distintos\dots{}
        \textbf{Importa el orden y entran todos}Permutaciones: \fmm{P_n=n!}. Ordenar personas en una fila, libros en una estantería, todas las letras de una palabra\dots{}
        \textbf{Entran todos, pero hay repetidos}Permutaciones con repetición: \fmm{PR_n^{\,a,b,\dots}=\frac{n!}{a!\,b!\cdots}}. Palabras con letras repetidas, secuencias con un número fijo de caras\dots{}
        \textbf{No importa el orden}Combinaciones: \fmm{C_{n,k}=\binom nk}. Grupos, equipos, comisiones, saludos, apuestas de lotería, sacar varias bolas o cartas a la vez\dots{}

\begin{fmnote}Un truco para saber si importa el orden: cambia el orden de dos elementos. Si obtienes una cosa distinta (otro podio, otro número), importa; si es lo mismo (el mismo equipo), no importa.\end{fmnote}
\end{fmprop}

\Needspace{5\baselineskip}
\begin{fmex}{Ejemplos resueltos}
\begin{enumerate}[label=\textbf{\arabic*.},leftmargin=*]
\item En un grupo de 10 personas hay que formar una comisión de 3. ¿De cuántas maneras se puede hacer? ¿Y si los tres tienen cargos distintos (presidente, secretario y tesorero)?

Una comisión es un grupo: no importa el orden. \fmm{C_{10,3}=\binom{10}{3}=\dfrac{10\cdot9\cdot8}{3\cdot2\cdot1}=120}.

Con cargos, sí importa el orden: \fmm{V_{10,3}=10\cdot9\cdot8=720} (cada comisión da \fmm{3!=6} repartos de cargos).
\item En la lotería Primitiva se marcan 6 números del 1 al 49. ¿Cuántas apuestas distintas hay?

No importa el orden en que se marcan: \fmm{\binom{49}{6}=\dfrac{49\cdot48\cdot47\cdot46\cdot45\cdot44}{6!}=13\,983\,816}.
\item Ocho amigos se saludan y cada uno da la mano a todos los demás una vez. ¿Cuántos apretones de manos hay?

Cada apretón es una pareja (A da la mano a B es lo mismo que B a A): \fmm{\binom82=\dfrac{8\cdot7}{2}=28}.
\item En una clase hay 12 chicas y 10 chicos. Hay que elegir una comisión de 2 chicas y 2 chicos. ¿De cuántas maneras se puede hacer?

Elegimos las chicas, \fmm{\binom{12}{2}=66}, y los chicos, \fmm{\binom{10}{2}=45}. Por el principio de multiplicación: \fmm{66\cdot45=2\,970}.
\item ¿Cuántas diagonales tiene un polígono convexo de 10 lados?

Cada pareja de vértices determina un segmento: \fmm{\binom{10}{2}=45}. De estos, 10 son lados: \fmm{45-10=35} diagonales.
\end{enumerate}
\end{fmex}

\Needspace{5.7cm}
\section{Experimentos aleatorios. Espacio muestral y sucesos}\hypertarget{sec5}{}
\begin{fmdef}{Experimento aleatorio}
Un experimento es \textbf{aleatorio} cuando, aunque lo repitamos en las mismas condiciones, \textbf{no podemos predecir} el resultado que obtendremos, aunque sí conocemos todos los resultados posibles. Por ejemplo: lanzar un dado, sacar una carta de una baraja, elegir al azar una persona de una población y preguntarle si tiene coche\dots{}

En cambio, un experimento es \textbf{determinista} si el resultado se puede predecir: dejar caer una pelota desde 2 m y medir el tiempo que tarda en llegar al suelo.
\end{fmdef}

\Needspace{5\baselineskip}
\begin{fmdef}{Espacio muestral y sucesos}
\begin{itemize}
\item El \textbf{espacio muestral}, \fmm{E} (o \fmm{\Omega}), es el conjunto de \textbf{todos los resultados posibles} del experimento.
\item Un \textbf{suceso} es cualquier subconjunto del espacio muestral. Se designa con letras mayúsculas: \fmm{A}, \fmm{B}, \fmm{C}\dots{} Decimos que el suceso \fmm{A} \textbf{ocurre} si el resultado del experimento es uno de los elementos de \fmm{A}.
\item \textbf{Suceso elemental}: formado por un solo resultado. \textbf{Suceso compuesto}: formado por dos o más resultados.
\item \textbf{Suceso seguro}: el espacio muestral \fmm{E}, que ocurre siempre. \textbf{Suceso imposible}: el conjunto vacío \fmm{\varnothing}, que no ocurre nunca.
\end{itemize}
\end{fmdef}

\Needspace{5\baselineskip}
\begin{fmex}{Ejemplos resueltos}
\begin{enumerate}[label=\textbf{\arabic*.},leftmargin=*]
\item Lanzamos un dado y observamos el resultado. Escribe el espacio muestral y los sucesos \fmm{A=}«salir par», \fmm{B=}«salir más de 4» y \fmm{C=}«salir 7».

\fmm{E=\{1,2,3,4,5,6\}}

\fmm{A=\{2,4,6\}}, \fmm{\ B=\{5,6\}}, \fmm{\ C=\varnothing} (suceso imposible).
\item Lanzamos tres monedas. Escribe el espacio muestral y el suceso \fmm{A=}«obtener al menos dos caras».

Con un diagrama de árbol (C = cara, + = cruz) obtenemos 8 resultados:

\fmm{E=\{CCC,\ CC{+},\ C{+}C,\ C{+}{+},\ {+}CC,\ {+}C{+},\ {+}{+}C,\ {+}{+}{+}\}}

\fmm{A=\{CCC,\ CC{+},\ C{+}C,\ {+}CC\}}: 4 resultados.
\item Elegimos al azar un alumno de un instituto y miramos si aprueba Matemáticas y si aprueba Historia. ¿Cuál es el espacio muestral?

Si \fmm{M=}«aprueba Matemáticas» y \fmm{H=}«aprueba Historia», hay 4 posibilidades: aprueba las dos, solo Matemáticas, solo Historia o ninguna. Este es el tipo de espacio muestral que encontraremos en muchos problemas de la PAU, donde \textbf{no} sabemos cuántos alumnos hay de cada tipo sino solo los porcentajes.
\end{enumerate}
\end{fmex}

\Needspace{5.7cm}
\section{Operaciones con sucesos}\hypertarget{sec6}{}
\begin{fmdef}{Unión, intersección, contrario y diferencia}
\begin{center}\renewcommand{\arraystretch}{1.5}\setlength{\parskip}{0pt}\small\begin{tabularx}{\linewidth}{>{\raggedright\arraybackslash}p{3.2cm}>{\raggedright\arraybackslash}X>{\raggedright\arraybackslash}X>{\raggedright\arraybackslash}X}\toprule
\textbf{Operación} & \textbf{Notación} & \textbf{Ocurre cuando\dots{}} & \textbf{Se lee}\\\midrule
Unión & $A\cup B$ & ocurre $A$ \textbf{o} $B$ (al menos uno de los dos) & «A o B»\\[3pt]\end{tabularx}\par\nointerlineskip\vskip0pt
\small\begin{tabularx}{\linewidth}{>{\raggedright\arraybackslash}p{3.2cm}>{\raggedright\arraybackslash}X>{\raggedright\arraybackslash}X>{\raggedright\arraybackslash}X}Intersección & $A\cap B$ & ocurren $A$ \textbf{y} $B$ a la vez & «A y B»\\[3pt]\end{tabularx}\par\nointerlineskip\vskip0pt
\small\begin{tabularx}{\linewidth}{>{\raggedright\arraybackslash}p{3.2cm}>{\raggedright\arraybackslash}X>{\raggedright\arraybackslash}X>{\raggedright\arraybackslash}X}Contrario & $\overline{A}$ o $A^c$ & \textbf{no} ocurre $A$ & «no A»\\[3pt]\end{tabularx}\par\nointerlineskip\vskip0pt
\small\begin{tabularx}{\linewidth}{>{\raggedright\arraybackslash}p{3.2cm}>{\raggedright\arraybackslash}X>{\raggedright\arraybackslash}X>{\raggedright\arraybackslash}X}Diferencia & $A-B=A\cap\overline{B}$ & ocurre $A$ pero no $B$ & «solo A»\\[3pt]\bottomrule\end{tabularx}\end{center}

\par\noindent \begin{minipage}[t]{0.31\linewidth}\centering {\centering \includegraphics[width=\linewidth,height=3.4cm,keepaspectratio]{figs/6146264b7fba.pdf}\par}\par\smallskip{\centering\footnotesize \fmm{A\cup B}\par}\end{minipage}\hfill\begin{minipage}[t]{0.31\linewidth}\centering {\centering \includegraphics[width=\linewidth,height=3.4cm,keepaspectratio]{figs/95428cfcdf2e.pdf}\par}\par\smallskip{\centering\footnotesize \fmm{A\cap B}\par}\end{minipage}\hfill\begin{minipage}[t]{0.31\linewidth}\centering {\centering \includegraphics[width=\linewidth,height=3.4cm,keepaspectratio]{figs/5982258f0389.pdf}\par}\par\smallskip{\centering\footnotesize \fmm{\overline{A}}\par}\end{minipage}\par\medskip
\par\noindent \begin{minipage}[t]{0.31\linewidth}\centering {\centering \includegraphics[width=\linewidth,height=3.4cm,keepaspectratio]{figs/2ac3dfb17807.pdf}\par}\par\smallskip{\centering\footnotesize \fmm{A-B=A\cap\overline{B}}\par}\end{minipage}\hfill\begin{minipage}[t]{0.31\linewidth}\centering {\centering \includegraphics[width=\linewidth,height=3.4cm,keepaspectratio]{figs/07f4d6bb45f0.pdf}\par}\par\smallskip{\centering\footnotesize Solo uno de los dos\par}\end{minipage}\hfill\begin{minipage}[t]{0.31\linewidth}\centering {\centering \includegraphics[width=\linewidth,height=3.4cm,keepaspectratio]{figs/b01a325b8366.pdf}\par}\par\smallskip{\centering\footnotesize \fmm{\overline{A}\cap\overline{B}}: ninguno\par}\end{minipage}\par\medskip

\begin{fmnote}\fmuline{Traducción del lenguaje}: «al menos uno» \fmm{\to A\cup B}; «ambos», «a la vez» \fmm{\to A\cap B}; «ninguno de los dos», «ni A ni B» \fmm{\to\overline{A}\cap\overline{B}}; «solo A» \fmm{\to A\cap\overline{B}}; «solo uno» \fmm{\to(A\cap\overline{B})\cup(\overline{A}\cap B)}; «como mucho uno», «no ambos» \fmm{\to\overline{A\cap B}}.\end{fmnote}
\end{fmdef}

\fmboxneedalt{\par\noindent \begin{minipage}[t]{0.48\linewidth}\vspace{0pt}\raggedright Dos sucesos son \textbf{incompatibles} (o disjuntos) si \textbf{no pueden ocurrir a la vez}: \fmm{A\cap B=\varnothing}. Si no, son \textbf{compatibles}.

Un suceso y su contrario siempre son incompatibles y, además, entre los dos llenan el espacio muestral: \fmm{A\cup\overline{A}=E}, \fmm{A\cap\overline{A}=\varnothing}.\end{minipage}\hfill\begin{minipage}[t]{0.48\linewidth}\vspace{0pt}\raggedright \par\medskip\begin{minipage}{\linewidth}{\centering \includegraphics[width=\linewidth,height=5.2cm,keepaspectratio]{figs/2d7b4c8ddaeb.pdf}\par}\par\smallskip{\centering\footnotesize Incompatibles: \fmm{A\cap B=\varnothing}\par}\end{minipage}\par\medskip\end{minipage}\par\medskip}
\begin{fmdef}{Sucesos incompatibles}
\fmtwoin{\par\noindent \begin{minipage}[t]{0.48\linewidth}\vspace{0pt}\raggedright Dos sucesos son \textbf{incompatibles} (o disjuntos) si \textbf{no pueden ocurrir a la vez}: \fmm{A\cap B=\varnothing}. Si no, son \textbf{compatibles}.

Un suceso y su contrario siempre son incompatibles y, además, entre los dos llenan el espacio muestral: \fmm{A\cup\overline{A}=E}, \fmm{A\cap\overline{A}=\varnothing}.\end{minipage}\hfill\begin{minipage}[t]{0.48\linewidth}\vspace{0pt}\raggedright \par\medskip\begin{minipage}{\linewidth}{\centering \includegraphics[width=\linewidth,height=5.2cm,keepaspectratio]{figs/2d7b4c8ddaeb.pdf}\par}\par\smallskip{\centering\footnotesize Incompatibles: \fmm{A\cap B=\varnothing}\par}\end{minipage}\par\medskip\end{minipage}\par\medskip}{

Dos sucesos son \textbf{incompatibles} (o disjuntos) si \textbf{no pueden ocurrir a la vez}: \fmm{A\cap B=\varnothing}. Si no, son \textbf{compatibles}.

Un suceso y su contrario siempre son incompatibles y, además, entre los dos llenan el espacio muestral: \fmm{A\cup\overline{A}=E}, \fmm{A\cap\overline{A}=\varnothing}.

\begin{center}\begin{minipage}{0.62\linewidth}\centering \par\medskip\begin{minipage}{\linewidth}{\centering \includegraphics[width=\linewidth,height=5.2cm,keepaspectratio]{figs/2d7b4c8ddaeb.pdf}\par}\par\smallskip{\centering\footnotesize Incompatibles: \fmm{A\cap B=\varnothing}\par}\end{minipage}\par\medskip\end{minipage}\end{center}
}
\end{fmdef}\global\fmstackfalse

\Needspace{5\baselineskip}
\begin{fmprop}{Leyes de De Morgan}

\fmdisp{\overline{A\cup B}=\overline{A}\cap\overline{B}\qquad\qquad\overline{A\cap B}=\overline{A}\cup\overline{B}}

«No ocurrir ninguno de los dos» es lo mismo que «no ocurrir A \textbf{y} no ocurrir B»; «no ocurrir los dos a la vez» es lo mismo que «no ocurrir A \textbf{o} no ocurrir B». Al hacer el contrario, la unión se convierte en intersección y al revés. Son muy útiles porque en la PAU a menudo dan \fmm{P(\overline{A}\cap\overline{B})} o \fmm{P(\overline{A}\cup\overline{B})}.

\par\noindent \begin{minipage}[t]{0.31\linewidth}\centering {\centering \includegraphics[width=\linewidth,height=3.4cm,keepaspectratio]{figs/22d577f7a3aa.pdf}\par}\par\smallskip{\centering\footnotesize \fmm{\overline{A\cup B}=\overline{A}\cap\overline{B}}\par}\end{minipage}\hfill\begin{minipage}[t]{0.31\linewidth}\centering {\centering \includegraphics[width=\linewidth,height=3.4cm,keepaspectratio]{figs/d35248ddc97d.pdf}\par}\par\smallskip{\centering\footnotesize \fmm{\overline{A\cap B}=\overline{A}\cup\overline{B}}\par}\end{minipage}\hfill\hspace*{0.330\linewidth}\par\medskip
\end{fmprop}

\Needspace{5\baselineskip}
\begin{fmex}{Ejemplos resueltos}
\begin{enumerate}[label=\textbf{\arabic*.},leftmargin=*]
\item En el lanzamiento de un dado, sean \fmm{A=\{2,4,6\}} y \fmm{B=\{4,5,6\}}. Calcula \fmm{A\cup B}, \fmm{A\cap B}, \fmm{\overline{A}}, \fmm{A-B} y comprueba la primera ley de De Morgan.

\fmm{A\cup B=\{2,4,5,6\}}, \fmm{\ A\cap B=\{4,6\}}, \fmm{\ \overline{A}=\{1,3,5\}}, \fmm{\ A-B=\{2\}}

\fmm{\overline{A\cup B}=\{1,3\}} y \fmm{\overline{A}\cap\overline{B}=\{1,3,5\}\cap\{1,2,3\}=\{1,3\}}
\item Sean \fmm{M=}«aprobar Matemáticas» y \fmm{H=}«aprobar Historia». Expresa con operaciones: (a) aprobar solo Historia; (b) aprobar como mucho una; (c) no aprobar ninguna.

(a) \fmm{H\cap\overline{M}\qquad} (b) \fmm{\overline{M\cap H}=\overline{M}\cup\overline{H}\qquad} (c) \fmm{\overline{M}\cap\overline{H}=\overline{M\cup H}}
\end{enumerate}
\end{fmex}

\Needspace{5.7cm}
\section{La probabilidad. Frecuencia relativa y axiomas}\hypertarget{sec7}{}
\begin{fmdef}{Frecuencia relativa y ley de los grandes números}
Si repetimos un experimento \fmm{n} veces y el suceso \fmm{A} ocurre \fmm{n_A} veces, la \textbf{frecuencia relativa} de \fmm{A} es \fmm{f_r(A)=\dfrac{n_A}{n}}.

\textbf{Ley de los grandes números}: cuando el número de repeticiones crece, la frecuencia relativa se estabiliza alrededor de un valor. Ese valor es la \textbf{probabilidad} del suceso, \fmm{P(A)}. Así se estiman muchas probabilidades en la vida real (seguros, encuestas, control de calidad).
\end{fmdef}

\Needspace{5\baselineskip}
\begin{fmdef}{Definición axiomática (Kolmogórov)}
Una \textbf{probabilidad} es una función que asigna a cada suceso \fmm{A} un número \fmm{P(A)} y que cumple:

\begin{enumerate}
\item \fmm{P(A)\geq0} para cualquier suceso \fmm{A}.
\item \fmm{P(E)=1}.
\item Si \fmm{A} y \fmm{B} son incompatibles, \fmm{P(A\cup B)=P(A)+P(B)}.
\end{enumerate}

\begin{fmnote}En consecuencia, \textbf{toda probabilidad está entre 0 y 1}: \fmm{0\leq P(A)\leq1}, con \fmm{P(\varnothing)=0}. Si en un ejercicio te da una probabilidad negativa o mayor que 1, ¡hay un error! Además, la suma de las probabilidades de todos los sucesos elementales es 1.\end{fmnote}
\end{fmdef}

\Needspace{5\baselineskip}
\begin{fmex}{Ejemplo resuelto}
\begin{enumerate}[label=\textbf{\arabic*.},leftmargin=*]
\item Un dado está trucado de manera que la probabilidad de obtener un 6 es \fmm{\frac12} y las otras cinco caras tienen la misma probabilidad \fmm{p}. Calcula \fmm{p} y la probabilidad de obtener un número par. \fmpau{PAU}

Las probabilidades de los sucesos elementales suman 1: \fmm{5p+\frac12=1\Rightarrow p=\frac{1}{10}}.

\fmm{P(\text{par})=P(2)+P(4)+P(6)=\frac1{10}+\frac1{10}+\frac12=\frac{7}{10}}.
\end{enumerate}
\end{fmex}

\Needspace{7.3cm}
\section{Regla de Laplace. Cálculo de probabilidades con combinatoria}\hypertarget{sec8}{}
\begin{fmprop}{Regla de Laplace}
Si el espacio muestral es finito y \textbf{todos los resultados son igual de probables} (equiprobables), la probabilidad de un suceso \fmm{A} es:

\fmdisp{P(A)=\frac{\text{número de casos favorables a }A}{\text{número de casos posibles}}}

Las palabras «al azar», «moneda equilibrada», «dado no trucado» o «bola elegida al azar» nos indican que podemos aplicarla. Cuidado: \textbf{no} se puede aplicar si los resultados no son equiprobables (un dado trucado, o la suma de dos dados si tomamos \fmm{E=\{2,\dots,12\}}).
\end{fmprop}

\Needspace{5\baselineskip}
\begin{fmex}{Ejemplos resueltos}
\begin{enumerate}[label=\textbf{\arabic*.},leftmargin=*]
\item Sacamos una carta de una baraja española de 40 cartas (4 palos de 10 cartas: 1--7, sota, caballo y rey). Probabilidad de (a) sacar un oro; (b) sacar una figura; (c) sacar un as o una copa.

(a) \fmm{P=\frac{10}{40}=\frac14\qquad} (b) hay 12 figuras: \fmm{P=\frac{12}{40}=\frac{3}{10}}

(c) 4 ases + 10 copas, pero el as de copas lo hemos contado dos veces: \fmm{4+10-1=13} casos favorables, \fmm{P=\frac{13}{40}}.
\item Lanzamos dos dados. Probabilidad de que la suma sea 7, y de que al menos uno de los dos sea un 6.

Hay \fmm{6\cdot6=36} parejas equiprobables. Suma 7: (1,6), (2,5), (3,4), (4,3), (5,2), (6,1) \fmm{\Rightarrow P=\frac{6}{36}=\frac16}.

Al menos un 6: 6 parejas con el primer dado igual a 6, 6 con el segundo y (6,6) contada dos veces: \fmm{\frac{11}{36}}. (Más rápido con el contrario: \fmm{1-\frac{25}{36}=\frac{11}{36}}.)
\end{enumerate}
\end{fmex}

\Needspace{5\baselineskip}
\begin{fmdef}{Recuento de casos: resumen de 9.1--9.4}
Cuando hay muchos casos, no los escribimos todos: los \textbf{contamos} con las técnicas de los apartados 9.1--9.4. Los casos posibles y los favorables se tienen que contar \textbf{de la misma manera} (los dos con orden o los dos sin orden).

\begin{center}\renewcommand{\arraystretch}{1.8}\setlength{\parskip}{0pt}\small\begin{tabularx}{\linewidth}{>{\raggedright\arraybackslash}p{7.2cm}>{\raggedright\arraybackslash}X}\toprule
\textbf{} & \textbf{Fórmula}\\\midrule
Variaciones (importa el orden, sin repetición) & $V_{n,k}=\dfrac{n!}{(n-k)!}$\\[3pt]\end{tabularx}\par\nointerlineskip\vskip0pt
\small\begin{tabularx}{\linewidth}{>{\raggedright\arraybackslash}p{7.2cm}>{\raggedright\arraybackslash}X}Variaciones con repetición & $VR_{n,k}=n^k$\\[3pt]\end{tabularx}\par\nointerlineskip\vskip0pt
\small\begin{tabularx}{\linewidth}{>{\raggedright\arraybackslash}p{7.2cm}>{\raggedright\arraybackslash}X}Permutaciones (ordenarlos todos) & $P_n=n!$\\[3pt]\end{tabularx}\par\nointerlineskip\vskip0pt
\small\begin{tabularx}{\linewidth}{>{\raggedright\arraybackslash}p{7.2cm}>{\raggedright\arraybackslash}X}Permutaciones con repetición & $PR_n^{a,b,\dots}=\dfrac{n!}{a!\,b!\cdots}$\\[3pt]\end{tabularx}\par\nointerlineskip\vskip0pt
\small\begin{tabularx}{\linewidth}{>{\raggedright\arraybackslash}p{7.2cm}>{\raggedright\arraybackslash}X}Combinaciones (no importa el orden) & $C_{n,k}=\dbinom nk=\dfrac{n!}{k!\,(n-k)!}$\\[3pt]\bottomrule\end{tabularx}\end{center}

\begin{fmnote}Sacar \fmm{k} bolas \textbf{a la vez} es lo mismo que sacarlas una tras otra sin reemplazamiento: podemos contar con combinaciones o multiplicar a lo largo de un árbol (9.13). Las dos maneras dan el mismo resultado.\end{fmnote}
\end{fmdef}

\Needspace{5\baselineskip}
\begin{fmex}{Ejemplos resueltos con recuento}
\begin{enumerate}[label=\textbf{\arabic*.},leftmargin=*]
\item En una clase hay 15 chicas y 10 chicos. Se eligen 3 personas al azar para una comisión. Probabilidad de que sean tres chicas, y de que haya al menos un chico.

Casos posibles: grupos de 3 de 25 (no importa el orden), \fmm{\binom{25}{3}=2300}. Casos favorables: \fmm{\binom{15}{3}=455}.

\fmm{P(3\ \text{chicas})=\frac{455}{2300}=\frac{91}{460}\approx0{,}1978} (con un árbol: \fmm{\frac{15}{25}\cdot\frac{14}{24}\cdot\frac{13}{23}}, lo mismo).

Al menos un chico es el contrario: \fmm{1-\frac{91}{460}=\frac{369}{460}\approx0{,}8022}.
\item Se elige al azar un código de 4 cifras (de 0000 a 9999). Probabilidad de que todas las cifras sean distintas.

\fmm{P=\dfrac{V_{10,4}}{VR_{10,4}}=\dfrac{10\cdot9\cdot8\cdot7}{10^4}=\dfrac{5040}{10000}=0{,}504}
\item Cinco amigos, entre ellos Ana y Biel, se sientan al azar en una fila de 5 sillas. Probabilidad de que Ana y Biel queden juntos.

Casos posibles: \fmm{5!=120}. Favorables: «pegamos» a Ana y Biel como si fueran una persona (\fmm{4!} órdenes) y ellos dos se pueden intercambiar (\fmm{2}): \fmm{2\cdot4!=48}.

\fmm{P=\frac{48}{120}=\frac25=0{,}4}
\end{enumerate}
\end{fmex}

\Needspace{5.7cm}
\section{Fórmulas de la probabilidad}\hypertarget{sec9}{}
\begin{fmprop}{Propiedades de la probabilidad}
\begin{center}\renewcommand{\arraystretch}{1.5}\setlength{\parskip}{0pt}\small\begin{tabularx}{\linewidth}{>{\raggedright\arraybackslash}p{5.2cm}>{\raggedright\arraybackslash}X}\toprule
\textbf{Propiedad} & \textbf{Fórmula}\\\midrule
Suceso contrario & $P(\overline{A})=1-P(A)$\\[3pt]\end{tabularx}\par\nointerlineskip\vskip0pt
\small\begin{tabularx}{\linewidth}{>{\raggedright\arraybackslash}p{5.2cm}>{\raggedright\arraybackslash}X}Unión (caso general) & $P(A\cup B)=P(A)+P(B)-P(A\cap B)$\\[3pt]\end{tabularx}\par\nointerlineskip\vskip0pt
\small\begin{tabularx}{\linewidth}{>{\raggedright\arraybackslash}p{5.2cm}>{\raggedright\arraybackslash}X}Unión de sucesos incompatibles & $P(A\cup B)=P(A)+P(B)$\\[3pt]\end{tabularx}\par\nointerlineskip\vskip0pt
\small\begin{tabularx}{\linewidth}{>{\raggedright\arraybackslash}p{5.2cm}>{\raggedright\arraybackslash}X}Diferencia («solo A») & $P(A\cap\overline{B})=P(A-B)=P(A)-P(A\cap B)$\\[3pt]\end{tabularx}\par\nointerlineskip\vskip0pt
\small\begin{tabularx}{\linewidth}{>{\raggedright\arraybackslash}p{5.2cm}>{\raggedright\arraybackslash}X}Ninguno de los dos (De Morgan) & $P(\overline{A}\cap\overline{B})=P(\overline{A\cup B})=1-P(A\cup B)$\\[3pt]\end{tabularx}\par\nointerlineskip\vskip0pt
\small\begin{tabularx}{\linewidth}{>{\raggedright\arraybackslash}p{5.2cm}>{\raggedright\arraybackslash}X}No ambos (De Morgan) & $P(\overline{A}\cup\overline{B})=P(\overline{A\cap B})=1-P(A\cap B)$\\[3pt]\end{tabularx}\par\nointerlineskip\vskip0pt
\small\begin{tabularx}{\linewidth}{>{\raggedright\arraybackslash}p{5.2cm}>{\raggedright\arraybackslash}X}Solo uno de los dos & $P(A\cup B)-P(A\cap B)=P(A)+P(B)-2P(A\cap B)$\\[3pt]\end{tabularx}\par\nointerlineskip\vskip0pt
\small\begin{tabularx}{\linewidth}{>{\raggedright\arraybackslash}p{5.2cm}>{\raggedright\arraybackslash}X}Si $A\subset B$ & $P(A)\leq P(B)$\\[3pt]\bottomrule\end{tabularx}\end{center}

¿Por qué se resta \fmm{P(A\cap B)} en la unión? Porque, al sumar \fmm{P(A)+P(B)}, la parte común (la intersección) la hemos contado \textbf{dos veces}.

Para tres sucesos: \fmm{P(A\cup B\cup C)=P(A)+P(B)+P(C)-P(A\cap B)-P(A\cap C)-P(B\cap C)+P(A\cap B\cap C)}.

\begin{fmnote}\fmuline{La idea clave}: el diagrama de Venn de dos sucesos tiene \textbf{cuatro regiones} (solo A, A y B, solo B, ninguno). Sus probabilidades suman 1, y \textbf{cualquier} probabilidad que nos pregunten es la suma de algunas de estas regiones. Si conocemos suficientes datos, las podemos calcular todas.\end{fmnote}
\end{fmprop}

\Needspace{5\baselineskip}
\begin{fmex}{Ejemplos resueltos}
\begin{enumerate}[label=\textbf{\arabic*.},leftmargin=*]
\item Sean \fmm{A} y \fmm{B} dos sucesos con \fmm{P(A)=0{,}6}, \fmm{P(B)=0{,}5} y \fmm{P(A\cap B)=0{,}4}. Calcula \fmm{P(A\cup B)}, \fmm{P(\overline{A}\cap\overline{B})}, \fmm{P(\overline{A}\cap B)} y \fmm{P(\overline{A}\cup\overline{B})}. \fmpau{PAU País Vasco 2020}

\fmm{P(A\cup B)=0{,}6+0{,}5-0{,}4=0{,}7}

De Morgan: \fmm{P(\overline{A}\cap\overline{B})=1-P(A\cup B)=0{,}3}

Solo B: \fmm{P(\overline{A}\cap B)=P(B)-P(A\cap B)=0{,}5-0{,}4=0{,}1}

De Morgan: \fmm{P(\overline{A}\cup\overline{B})=1-P(A\cap B)=0{,}6}
\item Sean \fmm{A} y \fmm{B} con \fmm{P(A\cup B)=0{,}8}, \fmm{P(\overline{A})=0{,}5} y \fmm{P(A\cap B)=0{,}3}. Calcula \fmm{P(B)} y \fmm{P(\overline{A}\cup\overline{B})}. \fmpau{PAU Baleares 2020}

\fmm{P(A)=1-0{,}5=0{,}5}. De la fórmula de la unión: \fmm{0{,}8=0{,}5+P(B)-0{,}3\Rightarrow P(B)=0{,}6}

\fmm{P(\overline{A}\cup\overline{B})=1-P(A\cap B)=0{,}7}
\item \textbf{Tres sucesos.} En una ciudad, el 30 \% lee el diario A, el 25 \% el B y el 20 \% el C; el 10 \% lee A y B, el 8 \% A y C, el 5 \% B y C, y el 3 \% lee los tres. Probabilidad de que una persona lea algún diario, de que no lea ninguno y de que solo lea A.

\fmm{P(A\cup B\cup C)=0{,}3+0{,}25+0{,}2-0{,}1-0{,}08-0{,}05+0{,}03=0{,}55}

Ninguno: \fmm{P(\overline{A}\cap\overline{B}\cap\overline{C})=1-0{,}55=0{,}45} (De Morgan también vale para tres sucesos).

Solo A: a \fmm{P(A)} le quitamos lo que comparte con B y con C, pero \fmm{A\cap B\cap C} lo hemos quitado dos veces: \fmm{0{,}3-0{,}1-0{,}08+0{,}03=0{,}15}.
\end{enumerate}
\end{fmex}

\fmheadboxalt{
\section{Cálculo de probabilidades: con el espacio muestral y sin él}\hypertarget{sec10}{}
}{\par\noindent \begin{minipage}[t]{0.488\linewidth}\vspace{0pt}\begin{tcolorbox}[enhanced,colback=fmsurface,colframe=fmmuted!35,arc=2.5mm,boxrule=0.6pt,left=2.5mm,right=2.5mm,top=1.5mm,bottom=1.5mm,before upper=\raggedright,equal height group=cmp2]{\color{fmnavy}\bfseries Conocemos el espacio muestral}\par

Dados, monedas, cartas, urnas\dots{} Podemos \textbf{contar} los casos. Calculamos las probabilidades con la \textbf{regla de Laplace} (casos favorables / casos posibles), a menudo ayudándonos de una tabla o de un árbol, y podemos comprobarlo con las fórmulas.\end{tcolorbox}\end{minipage}\hfill\begin{minipage}[t]{0.488\linewidth}\vspace{0pt}\begin{tcolorbox}[enhanced,colback=fmsurface,colframe=fmmuted!35,arc=2.5mm,boxrule=0.6pt,left=2.5mm,right=2.5mm,top=1.5mm,bottom=1.5mm,before upper=\raggedright,equal height group=cmp2]{\color{fmnavy}\bfseries No conocemos el espacio muestral}\par

Nos dan \textbf{porcentajes} o probabilidades («el 60 \% de los clientes\dots{}»). No podemos contar: tenemos que trabajar con las \textbf{fórmulas}, un \textbf{diagrama de Venn} con las cuatro regiones o una \textbf{tabla de contingencia}.\end{tcolorbox}\end{minipage}\par\medskip}
\begin{fmdef}{Dos situaciones diferentes}
\fmtwoin{\par\noindent \begin{minipage}[t]{0.488\linewidth}\vspace{0pt}\begin{tcolorbox}[enhanced,colback=fmsurface,colframe=fmmuted!35,arc=2.5mm,boxrule=0.6pt,left=2.5mm,right=2.5mm,top=1.5mm,bottom=1.5mm,before upper=\raggedright,equal height group=cmp1]{\color{fmnavy}\bfseries Conocemos el espacio muestral}\par

Dados, monedas, cartas, urnas\dots{} Podemos \textbf{contar} los casos. Calculamos las probabilidades con la \textbf{regla de Laplace} (casos favorables / casos posibles), a menudo ayudándonos de una tabla o de un árbol, y podemos comprobarlo con las fórmulas.\end{tcolorbox}\end{minipage}\hfill\begin{minipage}[t]{0.488\linewidth}\vspace{0pt}\begin{tcolorbox}[enhanced,colback=fmsurface,colframe=fmmuted!35,arc=2.5mm,boxrule=0.6pt,left=2.5mm,right=2.5mm,top=1.5mm,bottom=1.5mm,before upper=\raggedright,equal height group=cmp1]{\color{fmnavy}\bfseries No conocemos el espacio muestral}\par

Nos dan \textbf{porcentajes} o probabilidades («el 60 \% de los clientes\dots{}»). No podemos contar: tenemos que trabajar con las \textbf{fórmulas}, un \textbf{diagrama de Venn} con las cuatro regiones o una \textbf{tabla de contingencia}.\end{tcolorbox}\end{minipage}\par\medskip}{\begin{tcolorbox}[enhanced,breakable,colback=fmsurface,colframe=fmmuted!35,arc=2.5mm,boxrule=0.6pt,left=2.5mm,right=2.5mm,top=1.5mm,bottom=1.5mm,before upper=\raggedright]{\color{fmnavy}\bfseries Conocemos el espacio muestral}\par

Dados, monedas, cartas, urnas\dots{} Podemos \textbf{contar} los casos. Calculamos las probabilidades con la \textbf{regla de Laplace} (casos favorables / casos posibles), a menudo ayudándonos de una tabla o de un árbol, y podemos comprobarlo con las fórmulas.\end{tcolorbox}\begin{tcolorbox}[enhanced,breakable,colback=fmsurface,colframe=fmmuted!35,arc=2.5mm,boxrule=0.6pt,left=2.5mm,right=2.5mm,top=1.5mm,bottom=1.5mm,before upper=\raggedright]{\color{fmnavy}\bfseries No conocemos el espacio muestral}\par

Nos dan \textbf{porcentajes} o probabilidades («el 60 \% de los clientes\dots{}»). No podemos contar: tenemos que trabajar con las \textbf{fórmulas}, un \textbf{diagrama de Venn} con las cuatro regiones o una \textbf{tabla de contingencia}.\end{tcolorbox}}
\end{fmdef}\global\fmstackfalse

\Needspace{5\baselineskip}
\begin{fmdef}{Tabla de contingencia}
Es una tabla de doble entrada con dos sucesos y sus contrarios. En las celdas interiores están las \textbf{intersecciones}, y en los márgenes, las probabilidades de cada suceso. Cada fila y cada columna \textbf{suman} el total del margen, y el total general es 1 (o 100 \%). Con tres datos bien elegidos podemos rellenar toda la tabla.

\begin{center}\renewcommand{\arraystretch}{1.55}\setlength{\tabcolsep}{9pt}\begin{tabular}{l|cc|c}\toprule
\textbf{} & \textbf{$B$} & \textbf{$\overline{B}$} & \textbf{Total}\\\midrule
\textbf{$A$} & $P(A\cap B)$ & $P(A\cap\overline{B})$ & \textbf{$P(A)$}\\
\textbf{$\overline{A}$} & $P(\overline{A}\cap B)$ & $P(\overline{A}\cap\overline{B})$ & \textbf{$P(\overline{A})$}\\
\midrule\textbf{Total} & \textbf{$P(B)$} & \textbf{$P(\overline{B})$} & \textbf{1}\\
\bottomrule\end{tabular}\end{center}
\end{fmdef}

\Needspace{5\baselineskip}
\begin{fmex}{Ejemplos resueltos}
\begin{enumerate}[label=\textbf{\arabic*.},leftmargin=*]
\item \textbf{Conocemos el espacio muestral.} Lanzamos dos dados. Sean \fmm{A=}«la suma es 8» y \fmm{B=}«los dos dados son iguales». Calcula \fmm{P(A\cup B)}.

\fmm{A=\{(2,6),(3,5),(4,4),(5,3),(6,2)\}\Rightarrow P(A)=\frac{5}{36}}; \fmm{\ B=\{(1,1),\dots,(6,6)\}\Rightarrow P(B)=\frac{6}{36}}

\fmm{A\cap B=\{(4,4)\}\Rightarrow P(A\cap B)=\frac1{36}\Rightarrow P(A\cup B)=\frac{5+6-1}{36}=\frac{10}{36}=\frac{5}{18}}

Comprobación por recuento: \fmm{A\cup B} tiene \fmm{5+6-1=10} parejas.
\item \textbf{No conocemos el espacio muestral.} En un centro, el 50 \% del alumnado aprueba Historia, el 70 \% aprueba Matemáticas y el 30 \% aprueba ambas. Elegimos un alumno al azar. Calcula la probabilidad de que (a) apruebe solo una de las dos; (b) no apruebe más de una. \fmpau{PAU Andalucía 2025}

Rellenamos la tabla de contingencia (\fmm{H} = aprueba Historia, \fmm{M} = aprueba Matemáticas) a partir de los datos (en negrita) restando:

\begin{center}\renewcommand{\arraystretch}{1.55}\setlength{\tabcolsep}{9pt}\begin{tabular}{l|cc|c}\toprule
\textbf{} & \textbf{$M$} & \textbf{$\overline{M}$} & \textbf{Total}\\\midrule
\textbf{$H$} & \cellcolor{fmyellow!45}0,30 & 0,20 & \cellcolor{fmyellow!45}\textbf{0,50}\\
\textbf{$\overline{H}$} & 0,40 & 0,10 & \textbf{0,50}\\
\midrule\textbf{Total} & \cellcolor{fmyellow!45}\textbf{0,70} & \textbf{0,30} & \textbf{1}\\
\bottomrule\end{tabular}\end{center}

(a) \fmm{P(H\cap\overline{M})+P(\overline{H}\cap M)=0{,}2+0{,}4=0{,}6}

(b) «No más de una» es el contrario de «ambas»: \fmm{1-0{,}3=0{,}7}
\item El 69 \% de los habitantes de una ciudad ve series, el 35 \% películas y el 18 \% no ve ni series ni películas. Probabilidad de que un habitante vea series o películas, y de que vea series pero no películas. \fmpau{PAU Andalucía 2019}

De Morgan: \fmm{P(S\cup F)=1-P(\overline{S}\cap\overline{F})=1-0{,}18=0{,}82}

\fmm{P(S\cap F)=0{,}69+0{,}35-0{,}82=0{,}22\Rightarrow P(S\cap\overline{F})=0{,}69-0{,}22=0{,}47}
\end{enumerate}
\end{fmex}

\fmheadboxalt{
\section{Probabilidad condicionada}\hypertarget{sec11}{}
}{\par\noindent \begin{minipage}[t]{0.48\linewidth}\vspace{0pt}\raggedright La probabilidad de \fmm{A} \textbf{sabiendo que} ha ocurrido \fmm{B} (con \fmm{P(B)\gt0}) es:

\fmdisp{P(A\,|\,B)=\frac{P(A\cap B)}{P(B)}}

Se lee «probabilidad de \fmm{A} condicionada a \fmm{B}» (también se escribe \fmm{P(A/B)}). La información «ha ocurrido \fmm{B}» \textbf{reduce el espacio muestral} a \fmm{B}: ahora solo miramos qué parte de \fmm{B} es también \fmm{A}.\end{minipage}\hfill\begin{minipage}[t]{0.48\linewidth}\vspace{0pt}\raggedright \par\medskip\begin{minipage}{\linewidth}{\centering \includegraphics[width=\linewidth,height=5.2cm,keepaspectratio]{figs/2b9e6ed86fc0.pdf}\par}\par\smallskip{\centering\footnotesize El espacio muestral se reduce a \fmm{B}\par}\end{minipage}\par\medskip\end{minipage}\par\medskip}
\begin{fmdef}{Probabilidad condicionada}
\fmtwoin{\par\noindent \begin{minipage}[t]{0.48\linewidth}\vspace{0pt}\raggedright La probabilidad de \fmm{A} \textbf{sabiendo que} ha ocurrido \fmm{B} (con \fmm{P(B)\gt0}) es:

\fmdisp{P(A\,|\,B)=\frac{P(A\cap B)}{P(B)}}

Se lee «probabilidad de \fmm{A} condicionada a \fmm{B}» (también se escribe \fmm{P(A/B)}). La información «ha ocurrido \fmm{B}» \textbf{reduce el espacio muestral} a \fmm{B}: ahora solo miramos qué parte de \fmm{B} es también \fmm{A}.\end{minipage}\hfill\begin{minipage}[t]{0.48\linewidth}\vspace{0pt}\raggedright \par\medskip\begin{minipage}{\linewidth}{\centering \includegraphics[width=\linewidth,height=5.2cm,keepaspectratio]{figs/2b9e6ed86fc0.pdf}\par}\par\smallskip{\centering\footnotesize El espacio muestral se reduce a \fmm{B}\par}\end{minipage}\par\medskip\end{minipage}\par\medskip}{

La probabilidad de \fmm{A} \textbf{sabiendo que} ha ocurrido \fmm{B} (con \fmm{P(B)\gt0}) es:

\fmdisp{P(A\,|\,B)=\frac{P(A\cap B)}{P(B)}}

Se lee «probabilidad de \fmm{A} condicionada a \fmm{B}» (también se escribe \fmm{P(A/B)}). La información «ha ocurrido \fmm{B}» \textbf{reduce el espacio muestral} a \fmm{B}: ahora solo miramos qué parte de \fmm{B} es también \fmm{A}.

\begin{center}\begin{minipage}{0.62\linewidth}\centering \par\medskip\begin{minipage}{\linewidth}{\centering \includegraphics[width=\linewidth,height=5.2cm,keepaspectratio]{figs/2b9e6ed86fc0.pdf}\par}\par\smallskip{\centering\footnotesize El espacio muestral se reduce a \fmm{B}\par}\end{minipage}\par\medskip\end{minipage}\end{center}
}

\begin{fmnote}\fmuline{Cómo reconocerla}: «sabiendo que», «si», «de los que\dots{}», «entre los que\dots{}», «de las mujeres, qué porcentaje\dots{}». Cuidado: \fmm{P(A\,|\,B)\neq P(B\,|\,A)} en general. «El 40 \% de los que juegan a baloncesto juegan a tenis» es \fmm{P(T\,|\,B)=0{,}4}, no \fmm{P(T\cap B)}.\end{fmnote}
\end{fmdef}\global\fmstackfalse

\Needspace{5\baselineskip}
\begin{fmprop}{Regla del producto y propiedades}

\fmdisp{P(A\cap B)=P(B)\cdot P(A\,|\,B)=P(A)\cdot P(B\,|\,A)}

\fmdisp{P(\overline{A}\,|\,B)=1-P(A\,|\,B)\qquad\qquad P(A\cap B\cap C)=P(A)\cdot P(B\,|\,A)\cdot P(C\,|\,A\cap B)}

La condicionada se comporta como una probabilidad dentro de \fmm{B}: en particular, el contrario se calcula restando de 1 \textbf{con la misma condición}. Pero ¡\textbf{no} es cierto que \fmm{P(A\,|\,\overline{B})=1-P(A\,|\,B)}!
\end{fmprop}

\Needspace{5\baselineskip}
\begin{fmex}{Ejemplos resueltos}
\begin{enumerate}[label=\textbf{\arabic*.},leftmargin=*]
\item Sean \fmm{A} y \fmm{B} con \fmm{P(A)=0{,}6}, \fmm{P(B)=0{,}8} y \fmm{P(A\cap\overline{B})=0{,}1}. Calcula \fmm{P(A\,|\,\overline{B})} y \fmm{P(A\cup B)}. \fmpau{PAU Madrid 2019}

\fmm{P(A\,|\,\overline{B})=\dfrac{P(A\cap\overline{B})}{P(\overline{B})}=\dfrac{0{,}1}{0{,}2}=0{,}5}

\fmm{P(A\cup B)=P(B)+P(A\cap\overline{B})=0{,}8+0{,}1=0{,}9} (B más «solo A»)
\item En una población, el 45 \% son hombres. El 27 \% es hombre y lector de prensa deportiva, y un 38,5 \% es mujer y no lee. (a) De las mujeres, ¿qué porcentaje lee prensa deportiva? (b) De los lectores, ¿qué porcentaje son hombres? \fmpau{PAU Galicia 2021}

\begin{center}\renewcommand{\arraystretch}{1.55}\setlength{\tabcolsep}{9pt}\begin{tabular}{l|cc|c}\toprule
\textbf{} & \textbf{Lee (L)} & \textbf{$\overline{L}$} & \textbf{Total}\\\midrule
\textbf{Hombre (H)} & \cellcolor{fmyellow!45}0,27 & 0,18 & \cellcolor{fmyellow!45}\textbf{0,45}\\
\textbf{Mujer} & 0,165 & \cellcolor{fmyellow!45}0,385 & \textbf{0,55}\\
\midrule\textbf{Total} & \textbf{0,435} & \textbf{0,565} & \textbf{1}\\
\bottomrule\end{tabular}\end{center}

(a) \fmm{P(L\,|\,\text{mujer})=\dfrac{0{,}165}{0{,}55}=0{,}3\to} 30 \%

(b) \fmm{P(\text{hombre}\,|\,L)=\dfrac{0{,}27}{0{,}435}\approx0{,}6207\to} 62 \% aproximadamente
\end{enumerate}
\end{fmex}

\Needspace{5.7cm}
\section{Sucesos dependientes e independientes}\hypertarget{sec12}{}
\begin{fmdef}{Independencia de sucesos}
Dos sucesos son \textbf{independientes} si el hecho de que ocurra uno \textbf{no cambia} la probabilidad del otro: \fmm{P(A\,|\,B)=P(A)}. Esto es equivalente a:

\fmdisp{A \text{ y } B \text{ independientes}\iff P(A\cap B)=P(A)\cdot P(B)}

Para comprobarlo en la PAU, calculamos \fmm{P(A)\cdot P(B)} y lo comparamos con \fmm{P(A\cap B)}. Si coinciden, son independientes; si no, son \textbf{dependientes}.

\begin{itemize}
\item Si \fmm{A} y \fmm{B} son independientes, también lo son \fmm{A} y \fmm{\overline{B}}, \fmm{\overline{A}} y \fmm{B}, y \fmm{\overline{A}} y \fmm{\overline{B}}. Por ejemplo, \fmm{P(\overline{A}\cap\overline{B})=P(\overline{A})\cdot P(\overline{B})}.
\item En experimentos que se repiten sin influirse (lanzar dos dados, dos personas que tiran a canasta, extracciones \textbf{con} reemplazamiento) los sucesos de cada prueba son independientes y las probabilidades se multiplican.
\end{itemize}

\begin{fmnote}\fmuline{¡No confundas incompatibles con independientes!} Incompatibles: \fmm{P(A\cap B)=0}, no pueden pasar a la vez (si pasa uno, el otro \textbf{seguro} que no: dependen mucho el uno del otro). Independientes: \fmm{P(A\cap B)=P(A)P(B)}. Dos sucesos con probabilidad positiva \textbf{no} pueden ser a la vez incompatibles e independientes.\end{fmnote}
\end{fmdef}

\Needspace{5\baselineskip}
\begin{fmprop}{Tres o más sucesos independientes. «Al menos uno» en $n$ pruebas}
Tres sucesos \fmm{A}, \fmm{B} y \fmm{C} son \textbf{independientes} si lo son dos a dos y, además, \fmm{P(A\cap B\cap C)=P(A)\,P(B)\,P(C)}. En la práctica (pruebas que no se influyen) multiplicamos directamente.

\fmdisp{P(\text{al menos uno})=1-P(\text{ninguno})=1-P(\overline{A_1})\,P(\overline{A_2})\cdots P(\overline{A_n})}

Si las \fmm{n} pruebas tienen la misma probabilidad de éxito \fmm{p}, queda \fmm{1-(1-p)^n}. Para saber \textbf{cuántas pruebas} hacen falta para que esta probabilidad supere un valor, resolvemos la inecuación con logaritmos.
\end{fmprop}

\Needspace{5\baselineskip}
\begin{fmex}{Ejemplos resueltos}
\begin{enumerate}[label=\textbf{\arabic*.},leftmargin=*]
\item En un instituto, el 64 \% aprueba Matemáticas, el 72 \% aprueba Inglés y el 78 \% aprueba alguna de las dos. ¿Son independientes los sucesos? \fmpau{PAU Castilla-La Mancha 2024}

\fmm{P(M\cap A)=0{,}64+0{,}72-0{,}78=0{,}58}

\fmm{P(M)\cdot P(A)=0{,}64\cdot0{,}72=0{,}4608\neq0{,}58\Rightarrow} \textbf{no} son independientes.
\item Alex encesta un tiro libre con probabilidad 0,65 y Fran con probabilidad 0,48, de manera independiente. Si los dos tiran, probabilidad de que solo enceste Alex, y de que enceste al menos uno. \fmpau{PAU Murcia 2021}

\fmm{P(A\cap\overline{F})=0{,}65\cdot0{,}52=0{,}338}

Al menos uno = contrario de «ninguno»: \fmm{1-0{,}35\cdot0{,}52=1-0{,}182=0{,}818}
\item Tres estudiantes intentan resolver un problema de forma independiente; las probabilidades de que lo resuelvan son \fmm{\frac12}, \fmm{\frac13} y \fmm{\frac14}. Probabilidad de que el problema quede resuelto, y de que lo resuelva exactamente uno de los tres.

Resuelto = al menos uno lo resuelve: \fmm{1-\frac12\cdot\frac23\cdot\frac34=1-\frac14=\frac34}.

Exactamente uno: tres casos incompatibles, \fmm{\frac12\cdot\frac23\cdot\frac34+\frac12\cdot\frac13\cdot\frac34+\frac12\cdot\frac23\cdot\frac14=\frac{6+3+2}{24}=\frac{11}{24}}.
\end{enumerate}
\end{fmex}

\Needspace{7.3cm}
\section{Experiencias compuestas. Experiencias dependientes e independientes}\hypertarget{sec13}{}
\begin{fmdef}{Experiencias compuestas independientes y dependientes}
Una experiencia aleatoria \textbf{compuesta} está formada por dos o más experiencias simples (las \textbf{etapas}). Si \fmm{A_1} es un resultado de la primera etapa y \fmm{A_2} uno de la segunda:

\fmtwoin{\par\noindent \begin{minipage}[t]{0.488\linewidth}\vspace{0pt}\begin{tcolorbox}[enhanced,colback=fmsurface,colframe=fmmuted!35,arc=2.5mm,boxrule=0.6pt,left=2.5mm,right=2.5mm,top=1.5mm,bottom=1.5mm,before upper=\raggedright,equal height group=cmp3]{\color{fmnavy}\bfseries Experiencias independientes}\par

El resultado de una etapa \textbf{no influye} en las probabilidades de las demás: lanzar dos monedas, un dado y una moneda, extracciones \textbf{con} reemplazamiento\dots{}

\fmm{P(A_1\cap A_2)=P(A_1)\cdot P(A_2)}\end{tcolorbox}\end{minipage}\hfill\begin{minipage}[t]{0.488\linewidth}\vspace{0pt}\begin{tcolorbox}[enhanced,colback=fmsurface,colframe=fmmuted!35,arc=2.5mm,boxrule=0.6pt,left=2.5mm,right=2.5mm,top=1.5mm,bottom=1.5mm,before upper=\raggedright,equal height group=cmp3]{\color{fmnavy}\bfseries Experiencias dependientes}\par

El resultado de la primera etapa \textbf{cambia} las probabilidades de la segunda: extracciones \textbf{sin} reemplazamiento, elegir una urna y después una bola\dots{}

\fmm{P(A_1\cap A_2)=P(A_1)\cdot P(A_2\,|\,A_1)}\end{tcolorbox}\end{minipage}\par\medskip}{\begin{tcolorbox}[enhanced,breakable,colback=fmsurface,colframe=fmmuted!35,arc=2.5mm,boxrule=0.6pt,left=2.5mm,right=2.5mm,top=1.5mm,bottom=1.5mm,before upper=\raggedright]{\color{fmnavy}\bfseries Experiencias independientes}\par

El resultado de una etapa \textbf{no influye} en las probabilidades de las demás: lanzar dos monedas, un dado y una moneda, extracciones \textbf{con} reemplazamiento\dots{}

\fmm{P(A_1\cap A_2)=P(A_1)\cdot P(A_2)}\end{tcolorbox}\begin{tcolorbox}[enhanced,breakable,colback=fmsurface,colframe=fmmuted!35,arc=2.5mm,boxrule=0.6pt,left=2.5mm,right=2.5mm,top=1.5mm,bottom=1.5mm,before upper=\raggedright]{\color{fmnavy}\bfseries Experiencias dependientes}\par

El resultado de la primera etapa \textbf{cambia} las probabilidades de la segunda: extracciones \textbf{sin} reemplazamiento, elegir una urna y después una bola\dots{}

\fmm{P(A_1\cap A_2)=P(A_1)\cdot P(A_2\,|\,A_1)}\end{tcolorbox}}

\begin{fmnote}Para tres etapas: \fmm{P(A_1\cap A_2\cap A_3)=P(A_1)\,P(A_2\,|\,A_1)\,P(A_3\,|\,A_1\cap A_2)}. En los dos casos lo organizamos con un \textbf{diagrama de árbol}.\end{fmnote}
\end{fmdef}

\fmboxneedalt{\par\noindent \begin{minipage}[t]{0.48\linewidth}\vspace{0pt}\raggedright Una experiencia \textbf{compuesta} está formada por varias pruebas seguidas (sacar dos bolas, lanzar una moneda y después un dado\dots{}). La representamos con un \textbf{diagrama de árbol}:

\begin{itemize}
\item En cada rama escribimos la probabilidad \textbf{condicionada} a lo que ha pasado antes.
\item Las ramas que salen de un mismo nudo suman 1.
\item La probabilidad de un camino es el \textbf{producto} de las probabilidades de sus ramas (regla del producto).
\item Si un suceso corresponde a varios caminos, \textbf{sumamos} sus probabilidades.
\end{itemize}\end{minipage}\hfill\begin{minipage}[t]{0.48\linewidth}\vspace{0pt}\raggedright \par\medskip\begin{minipage}{\linewidth}{\centering \includegraphics[width=\linewidth,height=5.2cm,keepaspectratio]{figs/06577e522195.pdf}\par}\par\smallskip{\centering\footnotesize Urna con 3 R y 2 B, dos extracciones sin reemplazamiento\par}\end{minipage}\par\medskip\end{minipage}\par\medskip}
\begin{fmdef}{Diagramas de árbol}
\fmtwoin{\par\noindent \begin{minipage}[t]{0.48\linewidth}\vspace{0pt}\raggedright Una experiencia \textbf{compuesta} está formada por varias pruebas seguidas (sacar dos bolas, lanzar una moneda y después un dado\dots{}). La representamos con un \textbf{diagrama de árbol}:

\begin{itemize}
\item En cada rama escribimos la probabilidad \textbf{condicionada} a lo que ha pasado antes.
\item Las ramas que salen de un mismo nudo suman 1.
\item La probabilidad de un camino es el \textbf{producto} de las probabilidades de sus ramas (regla del producto).
\item Si un suceso corresponde a varios caminos, \textbf{sumamos} sus probabilidades.
\end{itemize}\end{minipage}\hfill\begin{minipage}[t]{0.48\linewidth}\vspace{0pt}\raggedright \par\medskip\begin{minipage}{\linewidth}{\centering \includegraphics[width=\linewidth,height=5.2cm,keepaspectratio]{figs/06577e522195.pdf}\par}\par\smallskip{\centering\footnotesize Urna con 3 R y 2 B, dos extracciones sin reemplazamiento\par}\end{minipage}\par\medskip\end{minipage}\par\medskip}{

Una experiencia \textbf{compuesta} está formada por varias pruebas seguidas (sacar dos bolas, lanzar una moneda y después un dado\dots{}). La representamos con un \textbf{diagrama de árbol}:

\begin{itemize}
\item En cada rama escribimos la probabilidad \textbf{condicionada} a lo que ha pasado antes.
\item Las ramas que salen de un mismo nudo suman 1.
\item La probabilidad de un camino es el \textbf{producto} de las probabilidades de sus ramas (regla del producto).
\item Si un suceso corresponde a varios caminos, \textbf{sumamos} sus probabilidades.
\end{itemize}

\begin{center}\begin{minipage}{0.62\linewidth}\centering \par\medskip\begin{minipage}{\linewidth}{\centering \includegraphics[width=\linewidth,height=5.2cm,keepaspectratio]{figs/06577e522195.pdf}\par}\par\smallskip{\centering\footnotesize Urna con 3 R y 2 B, dos extracciones sin reemplazamiento\par}\end{minipage}\par\medskip\end{minipage}\end{center}
}
\end{fmdef}\global\fmstackfalse

\fmboxneedalt{\par\noindent \begin{minipage}[t]{0.488\linewidth}\vspace{0pt}\begin{tcolorbox}[enhanced,colback=fmsurface,colframe=fmmuted!35,arc=2.5mm,boxrule=0.6pt,left=2.5mm,right=2.5mm,top=1.5mm,bottom=1.5mm,before upper=\raggedright,equal height group=cmp5]{\color{fmnavy}\bfseries Con reemplazamiento (con devolución)}\par

Devolvemos la bola a la urna antes de la siguiente extracción. La composición \textbf{no cambia}: las extracciones son \textbf{independientes} y las probabilidades de cada rama se repiten.

\fmm{P(R_1\cap R_2)=\frac35\cdot\frac35=\frac{9}{25}}\end{tcolorbox}\end{minipage}\hfill\begin{minipage}[t]{0.488\linewidth}\vspace{0pt}\begin{tcolorbox}[enhanced,colback=fmsurface,colframe=fmmuted!35,arc=2.5mm,boxrule=0.6pt,left=2.5mm,right=2.5mm,top=1.5mm,bottom=1.5mm,before upper=\raggedright,equal height group=cmp5]{\color{fmnavy}\bfseries Sin reemplazamiento (sin devolución)}\par

No devolvemos la bola. La composición \textbf{cambia} en cada extracción: las extracciones son \textbf{dependientes}. Sacarlas «a la vez» (dos bolas a la vez) es equivalente a sacarlas una tras otra sin reemplazamiento.

\fmm{P(R_1\cap R_2)=\frac35\cdot\frac24=\frac{3}{10}}\end{tcolorbox}\end{minipage}\par\medskip}
\begin{fmdef}{Con reemplazamiento y sin reemplazamiento}
\fmtwoin{\par\noindent \begin{minipage}[t]{0.488\linewidth}\vspace{0pt}\begin{tcolorbox}[enhanced,colback=fmsurface,colframe=fmmuted!35,arc=2.5mm,boxrule=0.6pt,left=2.5mm,right=2.5mm,top=1.5mm,bottom=1.5mm,before upper=\raggedright,equal height group=cmp4]{\color{fmnavy}\bfseries Con reemplazamiento (con devolución)}\par

Devolvemos la bola a la urna antes de la siguiente extracción. La composición \textbf{no cambia}: las extracciones son \textbf{independientes} y las probabilidades de cada rama se repiten.

\fmm{P(R_1\cap R_2)=\frac35\cdot\frac35=\frac{9}{25}}\end{tcolorbox}\end{minipage}\hfill\begin{minipage}[t]{0.488\linewidth}\vspace{0pt}\begin{tcolorbox}[enhanced,colback=fmsurface,colframe=fmmuted!35,arc=2.5mm,boxrule=0.6pt,left=2.5mm,right=2.5mm,top=1.5mm,bottom=1.5mm,before upper=\raggedright,equal height group=cmp4]{\color{fmnavy}\bfseries Sin reemplazamiento (sin devolución)}\par

No devolvemos la bola. La composición \textbf{cambia} en cada extracción: las extracciones son \textbf{dependientes}. Sacarlas «a la vez» (dos bolas a la vez) es equivalente a sacarlas una tras otra sin reemplazamiento.

\fmm{P(R_1\cap R_2)=\frac35\cdot\frac24=\frac{3}{10}}\end{tcolorbox}\end{minipage}\par\medskip}{\begin{tcolorbox}[enhanced,breakable,colback=fmsurface,colframe=fmmuted!35,arc=2.5mm,boxrule=0.6pt,left=2.5mm,right=2.5mm,top=1.5mm,bottom=1.5mm,before upper=\raggedright]{\color{fmnavy}\bfseries Con reemplazamiento (con devolución)}\par

Devolvemos la bola a la urna antes de la siguiente extracción. La composición \textbf{no cambia}: las extracciones son \textbf{independientes} y las probabilidades de cada rama se repiten.

\fmm{P(R_1\cap R_2)=\frac35\cdot\frac35=\frac{9}{25}}\end{tcolorbox}\begin{tcolorbox}[enhanced,breakable,colback=fmsurface,colframe=fmmuted!35,arc=2.5mm,boxrule=0.6pt,left=2.5mm,right=2.5mm,top=1.5mm,bottom=1.5mm,before upper=\raggedright]{\color{fmnavy}\bfseries Sin reemplazamiento (sin devolución)}\par

No devolvemos la bola. La composición \textbf{cambia} en cada extracción: las extracciones son \textbf{dependientes}. Sacarlas «a la vez» (dos bolas a la vez) es equivalente a sacarlas una tras otra sin reemplazamiento.

\fmm{P(R_1\cap R_2)=\frac35\cdot\frac24=\frac{3}{10}}\end{tcolorbox}}

\begin{fmnote}En los problemas aparecen también variantes: «se saca una bola y se reemplaza por dos del otro color», «se devuelve la bola y se añaden dos del mismo color»\dots{} El método es el mismo: en cada rama, la composición de la urna en ese momento.\end{fmnote}
\end{fmdef}\global\fmstackfalse

\Needspace{5\baselineskip}
\begin{fmex}{Ejemplos resueltos}
\begin{enumerate}[label=\textbf{\arabic*.},leftmargin=*]
\item Una urna contiene 5 bolas blancas, 3 negras y 2 rojas. Se sacan dos, consecutivamente y sin reemplazamiento. Calcula la probabilidad de que (a) las dos sean rojas; (b) sean del mismo color; (c) al menos una sea negra; (d) si la segunda es roja, que la primera haya sido blanca. \fmpau{PAU Cantabria 2023}

(a) \fmm{P(R_1\cap R_2)=\frac{2}{10}\cdot\frac19=\frac{1}{45}}

(b) \fmm{\frac{5}{10}\cdot\frac49+\frac{3}{10}\cdot\frac29+\frac{2}{10}\cdot\frac19=\frac{20+6+2}{90}=\frac{28}{90}=\frac{14}{45}}

(c) Contrario: ninguna negra. \fmm{1-\frac{7}{10}\cdot\frac69=1-\frac{42}{90}=\frac{8}{15}}

(d) \fmm{P(R_2)=\frac{5}{10}\cdot\frac29+\frac{3}{10}\cdot\frac29+\frac{2}{10}\cdot\frac19=\frac{18}{90}=\frac15}; \fmm{\ P(B_1\,|\,R_2)=\dfrac{P(B_1\cap R_2)}{P(R_2)}=\dfrac{10/90}{18/90}=\dfrac59}
\item Una urna tiene 5 bolas rojas y 3 verdes. Se saca una bola y se reemplaza por dos bolas del otro color. Después se saca una segunda bola. Probabilidad de que la segunda sea verde. \fmpau{PAU Andalucía 2017}

Si la primera es roja (\fmm{\frac58}), quedan 4 R y 5 V (9 bolas); si es verde (\fmm{\frac38}), quedan 7 R y 2 V (9 bolas).

\fmm{P(V_2)=\frac58\cdot\frac59+\frac38\cdot\frac29=\frac{25+6}{72}=\frac{31}{72}\approx0{,}4306}
\end{enumerate}
\end{fmex}

\Needspace{5.7cm}
\section{Teorema de la probabilidad total y teorema de Bayes}\hypertarget{sec14}{}
\begin{fmdef}{Sistema completo de sucesos}
Los sucesos \fmm{A_1, A_2,\dots,A_n} forman un \textbf{sistema completo} si son incompatibles dos a dos y su unión es todo el espacio muestral: siempre ocurre \textbf{uno y solo uno} de ellos. Por ejemplo: las fábricas que hacen un producto, las urnas de donde podemos sacar una bola, «estar enfermo / estar sano» (\fmm{A} y \fmm{\overline{A}}).
\end{fmdef}

\Needspace{5\baselineskip}
\begin{fmprop}{Teorema de la probabilidad total}

\fmdisp{P(B)=P(A_1)\,P(B\,|\,A_1)+P(A_2)\,P(B\,|\,A_2)+\dots+P(A_n)\,P(B\,|\,A_n)}

En el árbol: la probabilidad de \fmm{B} es la \textbf{suma de todos los caminos} que acaban en \fmm{B}. Responde preguntas como «¿cuál es la probabilidad de que una pieza sea defectuosa?» cuando cada fábrica tiene un porcentaje de defectos distinto.
\end{fmprop}

\Needspace{5\baselineskip}
\begin{fmprop}{Teorema de Bayes}

\fmdisp{P(A_i\,|\,B)=\frac{P(A_i)\,P(B\,|\,A_i)}{P(B)}=\frac{P(A_i)\,P(B\,|\,A_i)}{P(A_1)\,P(B\,|\,A_1)+\dots+P(A_n)\,P(B\,|\,A_n)}}

Bayes hace el camino \textbf{al revés}: sabemos el resultado \fmm{B} (la segunda prueba) y queremos la probabilidad de la causa \fmm{A_i} (la primera prueba). En el árbol: \textbf{el camino que nos interesa dividido entre la suma de todos los caminos que acaban en \fmm{B}}. Las probabilidades \fmm{P(A_i)} se llaman \emph{a priori} y las \fmm{P(A_i\,|\,B)}, \emph{a posteriori}.

\begin{fmnote}\fmuline{Pista}: si la pregunta es «sabiendo que la pieza es defectuosa, ¿qué probabilidad hay de que venga de la fábrica 2?», es decir, si nos dan la información de la \textbf{segunda} etapa y nos preguntan por la \textbf{primera}, es Bayes.\end{fmnote}
\end{fmprop}

\Needspace{5\baselineskip}
\begin{fmex}{Ejemplos resueltos}
\begin{enumerate}[label=\textbf{\arabic*.},leftmargin=*]
\item Tres profesores, A, B y C, han corregido el 25 \%, el 30 \% y el 45 \% de los exámenes de una oposición. Los porcentajes de errores de corrección son del 1 \%, el 2 \% y el 3 \%. (a) Probabilidad de que un examen esté mal corregido. (b) Si un examen está mal corregido, probabilidad de que lo haya corregido B. (c) ¿Qué corrector es el más probable? \fmpau{PAU Aragón 2021}

(a) Probabilidad total (\fmm{E} = error): \fmm{P(E)=0{,}25\cdot0{,}01+0{,}30\cdot0{,}02+0{,}45\cdot0{,}03=0{,}0025+0{,}006+0{,}0135=0{,}022}

(b) Bayes: \fmm{P(B\,|\,E)=\dfrac{0{,}006}{0{,}022}=\dfrac{3}{11}\approx0{,}2727}

(c) \fmm{P(A\,|\,E)=\frac{0{,}0025}{0{,}022}\approx0{,}1136}, \fmm{\ P(C\,|\,E)=\frac{0{,}0135}{0{,}022}\approx0{,}6136}. El más probable es C.
\item El 15 \% de una población tiene una enfermedad. Un test da positivo en el 92 \% de los enfermos y en el 4 \% de los sanos. Si el test de una persona da positivo, ¿qué probabilidad hay de que esté enferma? \fmpau{PAU Andalucía 2021}

\fmm{P(+)=0{,}15\cdot0{,}92+0{,}85\cdot0{,}04=0{,}138+0{,}034=0{,}172}

\fmm{P(M\,|\,+)=\dfrac{0{,}138}{0{,}172}=\dfrac{69}{86}\approx0{,}8023}
\end{enumerate}
\end{fmex}

  
\end{document}
